\chapter{Notes}
\label{ch:notes}

% ======================================================================
\section{Conventions and Units}
\label{sec:conventions_units}
% ======================================================================

This section collects the unit systems, notation, and Fourier
conventions used in this document.  The harmonic force
constants may originate either from phonopy finite-displacement
calculations or from Quantum ESPRESSO (QE) density-functional
perturbation theory, while the QuaTrEx solver uses its own internal
frequency-squared convention.  Since the later extraction, transport,
and anharmonic self-energy steps all rely on consistent units and phase
conventions, we summarize them here once and refer back to this section
throughout the chapter.

% ----------------------------------------------------------------------
\subsection{Units in phonopy and Quantum ESPRESSO}
\label{sec:phonopy_units}
% ----------------------------------------------------------------------

Phonopy stores force constants in eV/\AA$^2$, atomic masses in amu, and
atomic positions in \AA.  The dynamical matrix eigenvalues therefore
have units
\begin{equation}
  [\lambda] = \frac{\mathrm{eV}}{\mathrm{amu}\cdot\mathrm{\AA}^2}.
\end{equation}
Phonopy converts these eigenvalues to frequencies in THz via
\begin{equation}\label{eq:VaspToTHz}
  f\;[\mathrm{THz}]
  = \mathrm{sgn}(\lambda)\,\sqrt{|\lambda|}
    \times
    \frac{1}{2\pi}\sqrt{\frac{e}{m_u A^2}}  \approx 15.633\;\mathrm{THz},
\end{equation}
where $e = 1.602\times 10^{-19}$\,J is the elementary charge,
$m_u = 1.661\times 10^{-27}$\,kg is the atomic mass unit, and
$A = 10^{-10}$\,m is the metre-to-\AA\ conversion factor.

When phonopy is used together with the QE interface, the internal
representation inherited from QE uses atomic units: lattice vectors are
in bohr and force constants in Ry/bohr$^2$.  These quantities are
therefore converted before being used together with the rest of the
workflow.  The required relations are
\begin{align}
  a\;[\mathrm{\AA}]
    &= a\;[\mathrm{bohr}] \times a_0,
    & a_0 &= 0.52918\;\mathrm{\AA/bohr}, \\
  \Phi\;[\mathrm{eV/\AA^2}]
    &= \Phi\;[\mathrm{Ry/bohr^2}] \times \frac{E_h/2}{a_0^2},
    & \frac{E_h/2}{a_0^2} &\approx 48.59.
\end{align}
Here $E_h$ denotes the Hartree energy, so that $E_h/2 = 1$\,Ry.  In the
same way, Cartesian forces parsed from QE output are converted using
\begin{equation}
  1\;\mathrm{Ry/bohr} \approx 25.711\;\mathrm{eV/\AA}.
\end{equation}

The downstream workflow therefore treats eV, \AA, and amu as the common
intermediate unit system, independent of whether the original data came
from phonopy directly or through QE.  In other words, everything is
normalized to the phonopy eV/\AA/amu convention upon loading.

% ----------------------------------------------------------------------
\subsection{QuaTrEx solver units}
\label{sec:solver_units}
% ----------------------------------------------------------------------

The QuaTrEx phonon solver operates in angular frequency squared,
\begin{equation}
  \omega^2 \quad \text{in} \quad (\mathrm{rad/s})^2.
\end{equation}
Accordingly, the mass-weighted real-space blocks passed to the solver
must be converted from the phonopy eigenvalue units
eV/(amu$\cdot$\AA$^2$) to $(\mathrm{rad/s})^2$.  The corresponding
conversion factor is
\begin{equation}\label{eq:conversion}
  C = \frac{e}{m_u A^2}
  \approx 9.648 \times 10^{27}\;\mathrm{(rad/s)^2}.
\end{equation}
If $\Phi(0,\vc{R})$ denotes a harmonic force-constant block in
eV/\AA$^2$, then the real-space matrix block used by QuaTrEx is
\begin{equation}
  H(\vc{R})
  = C\,
    \frac{\Phi(0,\vc{R})}{\sqrt{m\,m'}},
\end{equation}
with $m$ and $m'$ in amu.  The resulting $H(\vc{R})$ therefore has units
of $(\mathrm{rad/s})^2$.

The same logic extends to cubic force constants.  After mass-weighting,
the third-order tensor
\begin{equation}
  \Psi_{\alpha\beta\gamma}
  = \frac{\Phi^{(3)}_{\alpha\beta\gamma}}
         {\sqrt{M_\kappa M_{\kappa'} M_{\kappa''}}}
\end{equation}
has units
\begin{equation}
  [\Psi] = \frac{\mathrm{eV}}
                 {\mathrm{\AA}^3\cdot\mathrm{amu}^{3/2}}.
\end{equation}
Its SI conversion factor is
\begin{equation}\label{eq:conversion_fc3}
  \texttt{CONVERSION\_FC3}
  = \frac{C}{A\sqrt{m_u}}
  = \frac{e}{m_u^{3/2} A^3}
  \approx 2.368\times10^{51},
\end{equation}
which converts
eV/(\AA$^3\cdot$amu$^{3/2}$) to the SI units required in the
phonon--phonon self-energy expressions.

% ----------------------------------------------------------------------
\subsection{Notation and Fourier conventions}
\label{sec:notation_fourier}
% ----------------------------------------------------------------------

We use bold symbols such as $\vc{R}$ and $\vc{q}$ for lattice vectors
and wavevectors.  Basis atoms within the unit cell are indexed by
Greek letters $\kappa,\kappa',\kappa''$, and Cartesian directions by
$\alpha,\beta,\gamma \in \{x,y,z\}$.  The displacement of atom
$\kappa$ in cell $l$ along Cartesian direction $\alpha$ is denoted
$u_\alpha(l\kappa)$.

The lattice vector of cell $l$ is written $\vc{R}_l$, while
$\vc{\tau}_\kappa$ denotes the basis-atom position within the unit
cell.  Unless stated otherwise, $\vc{\tau}_\kappa$ and $\vc{q}$ are
expressed in fractional crystal coordinates, so that factors of
$2\pi$ appear explicitly in all Fourier phases.

The harmonic force constants are written as
\begin{equation}
  \Phi_{\alpha\beta}(0\kappa,\,l'\kappa'),
\end{equation}
and the corresponding mass-weighted dynamical matrix at wavevector
$\vc{q}$ is a Fourier transform of these real-space couplings.

Two phase conventions are used in this work.  In the phonopy
convention,
\begin{equation}\label{eq:DA_conv}
  D_{\alpha\beta}^{(\mathrm{A})}(\kappa\kappa',\vc{q})
  = \frac{1}{\sqrt{m_\kappa m_{\kappa'}}}
    \sum_{l'} \Phi_{\alpha\beta}(0\kappa,\,l'\kappa')\,
    e^{2\pi i\,\vc{q}\cdot
      [\vc{R}_{l'} + \vc{\tau}_{\kappa'} - \vc{\tau}_\kappa]},
\end{equation}
whereas the QuaTrEx convention uses only the lattice vector in the
phase,
\begin{equation}\label{eq:DB_conv}
  D_{\alpha\beta}^{(\mathrm{B})}(\kappa\kappa',\vc{q})
  = \frac{1}{\sqrt{m_\kappa m_{\kappa'}}}
    \sum_{l'} \Phi_{\alpha\beta}(0\kappa,\,l'\kappa')\,
    e^{2\pi i\,\vc{q}\cdot\vc{R}_{l'}}.
\end{equation}
The two are related by a diagonal unitary gauge transformation,
\begin{equation}\label{eq:gauge_conv}
  D^{(\mathrm{B})}(\vc{q})
  = P(\vc{q})\,D^{(\mathrm{A})}(\vc{q})\,P^\dagger(\vc{q}),
\end{equation}
with
\begin{equation}
  P(\vc{q})_{\kappa\alpha,\kappa'\beta}
  = \delta_{\kappa\kappa'}\delta_{\alpha\beta}\,
    e^{2\pi i\,\vc{q}\cdot\vc{\tau}_\kappa}.
\end{equation}
Under Convention~B, the dynamical matrix is a pure lattice Fourier
series,
\begin{equation}\label{eq:fourier_series_conv}
  D^{(\mathrm{B})}(\vc{q})
  = \sum_{\vc{R}} H(\vc{R})\,e^{2\pi i\,\vc{q}\cdot\vc{R}}.
\end{equation}

For the transverse directions, we use the notation
$\vc{k}_\perp = (k_x,k_y)$ or $\vc{q}_\perp$ depending on context.  A
partial Fourier sum over transverse shifts then takes the form
\begin{equation}
  D(\vc{k}_\perp)
  = \sum_{(n_x,n_y)} M(n_x,n_y)\,
    e^{2\pi i (k_x n_x + k_y n_y)}.
\end{equation}
The transverse wavevectors are sampled on a Monkhorst--Pack mesh,
\begin{equation}
  k_i = \frac{n_i + 0.5}{N_i} - 0.5 + \mathrm{shift}_i,
  \qquad n_i = 0,\ldots,N_i-1.
\end{equation}

We also define the discrete Fourier-transform pair
\begin{align}
  G(\omega;\vq) &= \sum_{\vR} g(\omega;\vR)\,e^{-i\vq\cdot\vR},
  \label{eq:dft_fwd} \\
  g(\omega;\vR) &= \frac{1}{N}\sum_{\vq} G(\omega;\vq)\,e^{+i\vq\cdot\vR},
  \label{eq:dft_inv}
\end{align}
with the associated discrete convolution theorem
\begin{equation}\label{eq:conv_theorem}
  \frac{1}{N}\sum_{\vq'} A(\vq')\,B(\vq-\vq')
  = \sum_{\vR} a(\vR)\,b(\vR)\,e^{-i\vq\cdot\vR}.
\end{equation}

% ======================================================================
\section{Interatomic Force Constants}
\label{sec:ifc}
% ======================================================================

This section introduces the second- and third-order force
constants (FCs) used throughout the harmonic and anharmonic phonon
workflow.  FC2 and FC3 are generated as part of a common preprocessing
pipeline and stored in a unified output format.

% ----------------------------------------------------------------------
\subsection{Second- and third-order force constants}
\label{sec:ifc_definitions}
% ----------------------------------------------------------------------

In the harmonic approximation, the crystal potential energy is expanded
to second order in the atomic displacements from equilibrium.  The
second-order coefficients are the harmonic force constants,
\begin{equation}\label{eq:fc}
  \Phi_{\alpha\beta}(0\kappa,\,l'\kappa')
  = \frac{\partial^2 E}
         {\partial u_\alpha(0\kappa)\,\partial u_\beta(l'\kappa')},
\end{equation}
where $u_\alpha(l\kappa)$ denotes the displacement of atom $\kappa$ in
cell $l$ along Cartesian direction $\alpha$.  These coefficients relate
forces and displacements through
\begin{equation}\label{eq:force-disp}
  F_\alpha(0\kappa)
  = - \sum_{l'\kappa'\beta}
      \Phi_{\alpha\beta}(0\kappa,\,l'\kappa')\,
      u_\beta(l'\kappa').
\end{equation}
For a fixed atom pair $(0\kappa)$ and $(l'\kappa')$, the FCs form a
$3\times 3$ matrix, with rows corresponding to force components and
columns to displacement components.  If the unit cell contains $N$
atoms, then each lattice translation $\vc{R}_{l'}$ carries a
$(3N)\times(3N)$ force-constant block.

The corresponding dynamical matrix is obtained from the mass-weighted
Fourier transform of $\Phi$, and its eigenvalues and eigenvectors yield
the phonon frequencies and polarization vectors.

The third-order, or cubic, interatomic force constants are defined as
\begin{equation}\label{eq:fc3_def}
  \Phi^{(3)}_{\alpha\beta\gamma}(0\kappa,\,l'\kappa',\,l''\kappa'')
  = \frac{\partial^3 E}
         {\partial u_\alpha(0\kappa)\;
          \partial u_\beta(l'\kappa')\;
          \partial u_\gamma(l''\kappa'')}.
\end{equation}
By commutativity of partial derivatives, $\Phi^{(3)}$ is fully
symmetric under permutations of the three index triples
$(\alpha,0\kappa)$, $(\beta,l'\kappa')$, and
$(\gamma,l''\kappa'')$.

In the finite-displacement method, both FC2 and FC3 are obtained from
differences of DFT forces induced by small atomic displacements.  The
supercell size determines the real-space interaction range that can be
resolved before periodic wrap-around introduces aliasing.  For a
one-dimensional supercell of length $N$, the effective periodic force
constants satisfy
\begin{equation}
  \Phi_{\mathrm{per}}^{(N)}(r)
  = \sum_{p\in\mathbb Z} \Phi(r+pN).
\end{equation}
A larger supercell therefore increases the real-space range that can be
represented explicitly before distant interactions are folded back by
periodic boundary conditions.

% ----------------------------------------------------------------------
\subsection{Finite-displacement FCs}
\label{sec:finite_displacement}
% ----------------------------------------------------------------------

The FCs can be generated based on a finite-displacement workflow
using \texttt{phono3py}, QE, and the \texttt{symfc} fitting backend.  In this
approach, displaced supercells are generated explicitly, DFT forces are
computed for each displaced configuration, and both FC2 and FC3 are
reconstructed from the resulting force data.

The following steps are performed:
\begin{enumerate}
  \item \textbf{Structure setup.}
        The crystal structure is loaded and a phonopy object is
        created. The unit cell is loaded via \\
        \texttt{structure.py:load\_structure}, and a
        \texttt{Phonopy} object with a supercell is
        created by
        \texttt{structure.py:create\_phonopy\_from\_config}, which
        generates symmetry-inequivalent displacements.
  \item \textbf{DFT force calculations.}
        For each displaced supercell, a QE \texttt{pw.x}
        input is written and executed.  A QE \texttt{pw.x} SCF input
        file is written by\\
        \texttt{qe\_interface.py:write\_qe\_scf\_input}.  Existing
        converged outputs are detected and reused.  Forces parsed from
        QE output are converted from Ry/bohr to eV/\AA.

  \item \textbf{Force-constant reconstruction.}
        The set of forces is passed to phono3py, which reconstructs FC3
        and FC2 using the symmetry-adapted \texttt{symfc} backend, both
        FC2 and FC3 are written to a common \texttt{fc3.hdf5} file
        containing dense arrays in eV/\AA$^2$ and eV/\AA$^3$,
        respectively.
\end{enumerate}

For the Si FCC primitive cell with a $2\times2\times2$ supercell
(16 atoms), phono3py generates 114 symmetry-inequivalent displacements
for a pair cutoff of 5.5\,\AA.  On the mont-fort workstations, each QE SCF
calculation takes roughly 30\,s  with
\texttt{ecutwfc}=60\,Ry and a $2\times2\times2$ $k$-point mesh.

% ----------------------------------------------------------------------
\subsection{DFPT FCs}
\label{sec:dfpt}
% ----------------------------------------------------------------------

Third order FCs can be generated efficiently directly from density-functional
perturbation theory (DFPT) due to the $2n+1$ theorem.  One possibility is to obtain FC2 from QE's linear-response phonon code \texttt{ph.x} and compute FC3 with the D3Q plugin,
which implements third-order DFPT.  In contrast to the
finite-displacement method, the DFPT route works directly on the unit
cell and avoids supercell calculations. However, QE does only supports a restricted set of
norm-conserving potentials for third-order DFPT calculations, limiting the achievable accuracy.

The DFPT workflow does the following:
\begin{enumerate}
  \item \textbf{SCF ground state} (\texttt{pw.x}).
        A self-consistent calculation on the unit cell with a dense
        $k$-mesh (configured via \texttt{dfpt.kpoints}).

  \item \textbf{Phonon calculation} (\texttt{ph.x}).
        DFPT phonon calculation on a uniform $\vc{q}$-grid
        (\texttt{dfpt.q\_mesh}) with
        \texttt{ldisp\,=\,.true.} and \texttt{lqdir\,=\,.true.}.  This
        produces the dynamical matrices $D(\vc{q})$ and saves the
        first-order perturbation wavefunctions needed by D3Q.

  \item \textbf{FC2 Fourier transform} (\texttt{q2r.x}).
        The dynamical matrices are Fourier-transformed to real-space FC2
        via QE's \texttt{q2r.x}, producing the \texttt{flfrc} file.
        The acoustic sum rule is enforced
        (\texttt{zasr\,=\,'crystal'}).

  \item \textbf{Third-order DFPT} (\texttt{d3q.x}).
        For each $(\vc{q}_1,\vc{q}_2)$ pair on the grid, \texttt{d3q.x}
        computes the third-order dynamical tensor
        $D_3(\vc{q}_1,\vc{q}_2,-\vc{q}_1-\vc{q}_2)$.  The
        $\vc{q}$-points are specified in QE Cartesian coordinates
        ($2\pi/a_\text{lat}$ units), the conversion from crystal
        reciprocal coordinates is
        \begin{equation}\label{eq:q_conversion}
          \vc{q}_\text{cart}
          = q_1 \hat{\vc{b}}_1 + q_2 \hat{\vc{b}}_2 + q_3 \hat{\vc{b}}_3,
          \qquad
          \hat{B} = (A / |\vc{a}_1|)^{-\top},
        \end{equation}
        where $A$ is the matrix of lattice vectors (rows, in \AA) and
        $\hat{\vc{b}}_i$ are the rows of $\hat{B}$.

        For a $2\times2\times2$ grid there are $8^2 = 64$ pairs.
        Crystal symmetry can reduce this (e.g.\ to ${\sim}10$--15 for
        FCC Si with $O_h$ symmetry). The current implementation
        generates all pairs and relies on \texttt{d3q.x}'s internal
        symmetry handling.

  \item \textbf{FC3 Fourier transform} (\texttt{d3\_qq2rr.x}).
        The $(\vc{q}_1,\vc{q}_2)$-space third-order tensors are
        Fourier-transformed to real-space FC3
        $\Phi_3(\vc{R}_1,\vc{R}_2)$, producing the \texttt{mat3R} file.

  \item \textbf{Acoustic sum rule} (\texttt{d3\_asr.x}).
        Enforces translational invariance on the real-space FC3,
        $\sum_\kappa \Phi_3(\kappa,\kappa',\kappa'') = 0$ for each
        pair $(\kappa',\kappa'')$.  The ASR type is configurable
        (\texttt{asr}: \texttt{simple}, \texttt{crystal}, or
        \texttt{no}).

  \item \textbf{Sparsification} (\texttt{d3\_sparse.x}).
        Zeroes FC3 entries with magnitude below a threshold
        (\texttt{sparse\_thr}, in Ry/bohr$^3$), reducing storage
        and downstream computation cost.
\end{enumerate}

\paragraph{Output conversion.}
The \texttt{reap} step parses both the \texttt{q2r.x} FC2 output
(\texttt{fc2.dat}, in Ry/bohr$^2$) and the D3Q \texttt{mat3R} file (in
Ry/bohr$^3$) and converts to eV/\AA$^2$ and eV/\AA$^3$:
\begin{align}
  \Phi_2\;[\mathrm{eV/\AA^2}]
    &= \Phi_2\;[\mathrm{Ry/bohr^2}]
       \times \frac{E_h/2}{a_0^2},
    &&\approx 48.59, \label{eq:fc2_conv} \\
  \Phi_3\;[\mathrm{eV/\AA^3}]
    &= \Phi_3\;[\mathrm{Ry/bohr^3}]
       \times \frac{E_h/2}{a_0^3},
    &&\approx 91.84. \label{eq:fc3_conv}
\end{align}
The real-space lattice vectors $\vc{R}$ in the D3Q output (integer
crystal coordinates) are mapped to supercell atom indices
$j = (l_1 n_2 n_3 + l_2 n_3 + l_3)\, n_\text{at} + \kappa$ with
$l_i = R_i \bmod n_i$, matching phonopy's supercell convention.  FC2
is expanded to the full $(N_\text{super})^2$ array via translational
invariance.  Both arrays are saved to \texttt{fc3.hdf5} in the same
format as the finite-displacement method.

The DFPT-specific parameters are in the \texttt{dfpt} section of the
YAML config:
\begin{center}
\begin{tabular}{lll}
  \toprule
  Key & Default & Description \\
  \midrule
  \texttt{q\_mesh}   & $[2,2,2]$ & $\vc{q}$-grid for \texttt{ph.x} and \texttt{d3q.x} \\
  \texttt{kpoints}   & $[8,8,8]$ & SCF $k$-mesh for the unit cell \\
  \texttt{tr2\_ph}   & $10^{-14}$& \texttt{ph.x} self-consistency threshold \\
  \texttt{ph\_command}       & \texttt{ph.x}          & \texttt{ph.x} executable \\
  \texttt{d3q\_command}      & \texttt{d3q.x}         & \texttt{d3q.x} executable \\
  \texttt{d3\_qq2rr\_command}& \texttt{d3\_qq2rr.x}   & Fourier transform executable \\
  \texttt{d3\_asr\_command}  & \texttt{d3\_asr.x}     & Acoustic sum rule executable \\
  \texttt{d3\_sparse\_command}& \texttt{d3\_sparse.x} & Sparsification executable \\
  \texttt{asr}               & \texttt{simple}        & ASR type: \texttt{simple}, \texttt{crystal}, or \texttt{no} \\
  \texttt{sparse\_thr}       & $10^{-5}$              & FC3 threshold (Ry/bohr$^3$); \texttt{null} to skip \\
  \texttt{q2r\_command}      & \texttt{q2r.x}         & \texttt{q2r.x} executable \\
  \texttt{work\_dir}         & \texttt{./dfpt}        & Working directory \\
  \texttt{ph\_timeout}       & 7200\,s                & Timeout per \texttt{ph.x} run \\
  \texttt{d3q\_timeout}      & 14400\,s               & Timeout per \texttt{d3q.x} triplet \\
  \bottomrule
\end{tabular}
\end{center}

% ----------------------------------------------------------------------
\subsection{Mass-weighting and unit conversion}
\label{sec:ifc_mass_weighting}
% ----------------------------------------------------------------------

The harmonic transport problem is formulated in terms of mass-weighted
FC2 blocks.  Given a harmonic IFC block
$\Phi_{\alpha\beta}(0\kappa,l'\kappa')$, the corresponding mass-weighted
quantity is
\begin{equation}
  H_{\alpha\beta}(0\kappa,l'\kappa')
  = \frac{\Phi_{\alpha\beta}(0\kappa,l'\kappa')}
         {\sqrt{M_\kappa M_{\kappa'}}},
\end{equation}
with masses in amu.  In the common phonopy unit system,
$H_{\alpha\beta}$ therefore has units
eV/(amu$\cdot$\AA$^2$), which are later converted to
$(\mathrm{rad/s})^2$ using the factor defined in
\cref{sec:solver_units}.

For FC3, the analogous mass-weighted tensor is
\begin{equation}\label{eq:fc3_mw}
  \Psi_{\alpha\beta\gamma}(\kappa\kappa'\kappa'')
  = \frac{\Phi^{(3)}_{\alpha\beta\gamma}(\kappa\kappa'\kappa'')}
         {\sqrt{M_\kappa M_{\kappa'} M_{\kappa''}}},
\end{equation}
with units
\begin{equation}
  [\Psi]
  = \frac{\mathrm{eV}}
         {\mathrm{\AA}^3\cdot\mathrm{amu}^{3/2}}.
\end{equation}

For DFPT-generated FCs, the real-space force constants are parsed in
atomic units and converted as
\begin{align}
  \Phi_2\;[\mathrm{eV/\AA^2}]
    &= \Phi_2\;[\mathrm{Ry/bohr^2}]
       \times \frac{E_h/2}{a_0^2},
    &&\approx 48.59, \\
  \Phi_3\;[\mathrm{eV/\AA^3}]
    &= \Phi_3\;[\mathrm{Ry/bohr^3}]
       \times \frac{E_h/2}{a_0^3},
    &&\approx 91.84.
\end{align}

Finally, the FC3 tensor can be converted to SI units for use in the
anharmonic self-energy.  Using the harmonic conversion factor
$C = e/(m_u A^2)$, the cubic conversion becomes
\begin{equation}\label{eq:conversion_fc3_ifc}
  \texttt{CONVERSION\_FC3}
  = \frac{C}{A\sqrt{m_u}}
  = \frac{e}{m_u^{3/2} A^3}.
\end{equation}
This converts
eV/(\AA$^3\cdot$amu$^{3/2}$) into the corresponding SI units required
by the phonon--phonon self-energy expressions.

% ======================================================================
\section{QuaTrEx Input Generation}
\label{sec:quatrex_harmonic_input}
% ======================================================================

This section describes how the harmonic second-order force constants
(FC2) are converted into the real-space block representation required
by the QuaTrEx phonon solver.  In particular, we construct the
harmonic dynamical matrix, transform it from phonopy's phase
convention to the lattice-Fourier convention used by QuaTrEx,
extract real-space blocks by an inverse discrete Fourier transform,
and writing the resulting data to the solver input files. This can likely be
done more efficiently by avoiding the Full Dynamical matrix in Finite Displacement
or directly starting from the Dynamical matrix with DFPT.

% ----------------------------------------------------------------------
\subsection{Dynamical matrix and phase conventions}
\label{sec:dynmat}
% ----------------------------------------------------------------------

The dynamical matrix is the mass-weighted Fourier transform of
the second-order force constants.  Its eigenvalues are the squared
phonon frequencies and its eigenvectors the polarization vectors. There are
two different common conventions for the dynamical matrix:
phonopy includes basis-atom phase factors, QuaTrEx uses a pure lattice Fourier
series in the cell translations.

\subsubsection{Convention A (phonopy)}

In phonopy, the dynamical matrix at wavevector $\vc{q}$ is written as
\begin{equation}\label{eq:DA}
  D_{\alpha\beta}^{(\mathrm{A})}(\kappa\kappa',\vc{q})
  = \frac{1}{\sqrt{m_\kappa m_{\kappa'}}}
    \sum_{l'}
    \Phi_{\alpha\beta}(0\kappa,\,l'\kappa')\;
    e^{2\pi i\,\vc{q}\cdot
      [\vc{R}_{l'} + \vc{\tau}_{\kappa'} - \vc{\tau}_\kappa]},
\end{equation}
where $\vc{q}$ is expressed in fractional reciprocal coordinates,
$\vc{R}_{l'}$ is the lattice vector of cell $l'$, and
$\vc{\tau}_\kappa$ denotes the basis-atom position within the unit
cell.  The phase therefore depends on the full interatomic separation
\begin{equation}
  \vc{R}_{l'} + \vc{\tau}_{\kappa'} - \vc{\tau}_\kappa.
\end{equation}
This is what is returned by \texttt{phonon.get\_dynamical\_matrix\_at\_q(q)}.

\subsubsection{Convention B (QuaTrEx)}

QuaTrEx uses a lattice-only Fourier convention in which the phase
depends only on the cell translation:
\begin{equation}\label{eq:DB}
  D_{\alpha\beta}^{(\mathrm{B})}(\kappa\kappa',\vc{q})
  = \frac{1}{\sqrt{m_\kappa m_{\kappa'}}}
    \sum_{l'}
    \Phi_{\alpha\beta}(0\kappa,\,l'\kappa')\;
    e^{2\pi i\,\vc{q}\cdot \vc{R}_{l'}}.
\end{equation}
In this convention the dependence on basis-atom positions is removed
from the Fourier phase.  The dynamical matrix is then a pure lattice
Fourier series,
\begin{equation}\label{eq:fourier_series}
  D^{(\mathrm{B})}(\vc{q})
  = \sum_{\vc{R}} H(\vc{R})\;e^{2\pi i\,\vc{q}\cdot\vc{R}},
\end{equation}
where
\begin{equation}
  H(\vc{R}) = \frac{\Phi(0,\vc{R})}{\sqrt{m\,m'}}
\end{equation}
denotes the corresponding mass-weighted real-space block.

\subsubsection{Gauge transformation between conventions}
\label{sec:gauge}

The two conventions are related by a unitary similarity
transformation,
\begin{equation}\label{eq:gauge_eq}
  D^{(\mathrm{B})}(\vc{q})
  = P(\vc{q})\;D^{(\mathrm{A})}(\vc{q})\;P^\dagger(\vc{q}),
\end{equation}
with diagonal gauge matrix
\begin{equation}\label{eq:P}
  P(\vc{q})_{\kappa\alpha,\kappa'\beta}
  = \delta_{\kappa\kappa'}\delta_{\alpha\beta}\;
    e^{2\pi i\,\vc{q}\cdot\vc{\tau}_\kappa}.
\end{equation}
Each basis-atom phase is repeated three times, once for each Cartesian
degree of freedom. This is a similarity transformation and hence leaves
the eigenvalues of the dynamical matrix unchanged and only modifies the
phase convention of the eigenvectors.

% ----------------------------------------------------------------------
\subsection{Real-space block extraction}
\label{sec:blocks}
% ----------------------------------------------------------------------

The mass-weighted real-space blocks are obtained by an inverse discrete
Fourier transform on a uniform $\vc{q}$-mesh.  For a cubic
$N_q\times N_q\times N_q$ mesh, the extracted block at lattice
translation $(n_x,n_y,n_z)$ is
\begin{equation}\label{eq:idft}
  H(n_x,n_y,n_z)
  = \frac{C}{N_q^3}
    \sum_{i,j,k=0}^{N_q-1}
    D^{(\mathrm{B})}\!\left(\frac{i}{N_q},
                            \frac{j}{N_q},
                            \frac{k}{N_q}\right)
    e^{-2\pi i(n_x i + n_y j + n_z k)/N_q},
\end{equation}
where $C$ is the unit-conversion factor from
eV/(amu$\cdot$\AA$^2$) to $(\mathrm{rad/s})^2$, defined in
\cref{sec:solver_units}.

Because $D^{(\mathrm{B})}$ is a lattice Fourier series of real
force-constant matrices, the extracted blocks $H(\vc{R})$ are real up
to numerical roundoff.  Their imaginary part is therefore discarded. The blocks satisfy the symmetry relation
\begin{equation}\label{eq:sym}
  H(\vc{R})^T = H(-\vc{R}),
\end{equation}
which guarantees Hermiticity of the reconstructed dynamical matrix at
every $\vc{q}$.

% ----------------------------------------------------------------------
\subsection{QuaTrEx Input format}
\label{sec:quatrex_input}
% ----------------------------------------------------------------------

The QuaTrEx solver reads three input files which
specify the real-space Hamiltonian, the crystal structure, and the
transport parameters.

\subsubsection{Input files}
\label{sec:input_files}

\paragraph{\texttt{dynamical\_matrix.mat}}
This MATLAB \texttt{.mat} file stores the real-space blocks
$H(\vc{R})$ in Convention~B, already converted to
$(\mathrm{rad/s})^2$.  Each block is stored under a string key of the
form
\begin{equation}
  \texttt{"[nx, ny, nz]"},
\end{equation}
mapping to a real matrix of dimension
$(n_\mathrm{dof}\times n_\mathrm{dof})$, where
$n_\mathrm{dof}=3N_b$ and $N_b$ is the number of atoms in the unit
cell.  For nearest-neighbour interactions extracted on a
$3\times3\times3$ mesh, the file contains up to 27 blocks.

The solver interprets these blocks as real-space couplings and, for
periodic transverse directions, performs the Fourier sum
\begin{equation}
  D(\vc{k}_\perp)
  = \sum_{(n_x,n_y)} M(n_x,n_y)\;
    e^{2\pi i\,(k_x n_x + k_y n_y)}.
\end{equation}

\paragraph{\texttt{structure.xyz}}
This is an extended XYZ file containing the lattice vectors in the
header and Cartesian atomic positions in the body.  It follows the ASE
extended-XYZ convention, for example
\begin{verbatim}
8
Lattice="5.4018 0.0 0.0 0.0 5.4018 0.0 0.0 0.0 5.4018" ...
Si  0.00000000  0.00000000  0.00000000
...
\end{verbatim}
The loader (\texttt{distributed\_read\_xyz}) extracts the $3\times3$
lattice matrix and Cartesian atom positions.  The number of atoms times
\texttt{num\_orbitals\_per\_atom} determines $n_\mathrm{dof}$.

\paragraph{\texttt{quatrex\_config.toml}}
This file specifies the transport setup, such as transport direction,
number of transport cells, neighbour cutoff, transverse $k$-mesh,
broadening, and the open-boundary-condition algorithm.  Important fields are
\begin{center}
\begin{tabular}{lll}
  \toprule
  Field & Example value & Purpose \\
  \midrule
  \texttt{transport\_direction}       & \texttt{'z'}            & transport axis \\
  \texttt{construct\_from\_unit\_cell}& \texttt{true}           & input mode selector \\
  \texttt{num\_transport\_cells}      & 4                       & device length in unit cells \\
  \texttt{neighbor\_cell\_cutoff}     & \texttt{[1,1,1]}        & block-range cutoff \\
  \texttt{kpoint\_grid}               & \texttt{[4,4,1]}        & transverse mesh \\
  \texttt{kpoint\_shift}              & \texttt{[0,0,0]}        & Monkhorst--Pack shift \\
  \texttt{num\_orbitals\_per\_atom.Si}& 3                       & DOFs per atom \\
  \texttt{phonon.eta}                 & \texttt{1e-8}           & broadening $(\mathrm{rad/s})^2$ \\
  \texttt{phonon.obc.algorithm}       & \texttt{"sancho-rubio"} & OBC method \\
  \bottomrule
\end{tabular}
\end{center}

\subsubsection{Input modes}
\label{sec:input_modes}

QuaTrEx supports two input modes, selected by
\texttt{construct\_from\_unit\_cell}.  In both modes the input is a
MATLAB \texttt{.mat} file loaded by \texttt{scipy.io.loadmat}.  The
loader (\\\texttt{distributed\_load} in
\texttt{qttools/utils/mpi\_utils.py}) parses keys that start with
\texttt{[} and interprets them as integer tuples:
\begin{lstlisting}
arr = loadmat(path)
arr = {
    tuple(map(int, r.strip("[]").split(","))): h_r
    for r, h_r in arr.items()
    if r.startswith("[")
}
\end{lstlisting}
That is, a key string such as \texttt{"[0, 1, -1]"} becomes the Python
tuple \texttt{(0, 1, -1)}.  Each value is a 2D NumPy array (real or
complex) representing a matrix block.

\paragraph{Mode 1: pre-assembled device matrix}
If \texttt{construct\_from\_unit\_cell = false}, the
\texttt{dynamical\_matrix.mat} file contains a single pre-assembled
device matrix stored under the key \texttt{(0,0,0)}:
\begin{lstlisting}
matrix_sparray = distributed_load(f"{matrix_name}.mat")
assert (0, 0, 0) in matrix_sparray.keys()
matrix_sparray = matrix_sparray[(0, 0, 0)]
\end{lstlisting}
The result is a single sparse matrix of size
$(N_\mathrm{blocks}\times B)^2$, where $B$ is the block size and
$N_\mathrm{blocks}$ is the number of transport cells.  No transverse
$k$-point Fourier assembly is performed.

\paragraph{Mode 2: unit-cell blocks}
If \texttt{construct\_from\_unit\_cell = true}, the file contains the
dictionary of real-space unit-cell blocks
$H(n_x,n_y,n_z)$.  The solver loads all blocks, constructs a dense
translation grid, optionally trims it using
\texttt{neighbor\_cell\_cutoff}, expands the blocks into
transport-direction BTD matrices, and performs the transverse
Fourier assembly on the chosen Monkhorst--Pack mesh.

The function \texttt{\_load\_matrix\_from\_unit\_cell} loads all keys,
determines the bounding box from the minimum and maximum key
coordinates, and packs the blocks into a 5D array:
\begin{lstlisting}
# keys: array of all (nx, ny, nz) tuples
min_coords = keys.min(axis=0)
max_coords = keys.max(axis=0)
grid_shape = max_coords - min_coords + 1
unit_cells = zeros(grid_shape + matrix_shape)
for coord, matrix in matrices.items():
    unit_cells[coord] = matrix
\end{lstlisting}
The resulting array has shape
$(G_x,G_y,G_z,n_\mathrm{dof},n_\mathrm{dof})$, where
$(G_x,G_y,G_z)$ is the grid shape (e.g.\ $(3,3,3)$ for
$\{-1,0,+1\}^3$).  The loader checks that
$G_x \times G_y \times G_z = \text{number of keys}$, i.e.\ the grid
must be dense: no missing entries are permitted.  Missing blocks must
be stored as zero matrices.

If \texttt{neighbor\_cell\_cutoff} is set (e.g.\ \texttt{[1,1,1]}),
\texttt{trim\_tight\_binding\_matrix} zeros out all blocks beyond that
range.  For cutoff $(c_x,c_y,c_z)$, a block at index
$(n_x,n_y,n_z)$ (centred at zero) is kept only if
$|n_x| \le c_x$, $|n_y| \le c_y$, $|n_z| \le c_z$.

The function \texttt{expand\_tight\_binding\_matrix} converts the
unit-cell blocks into a BTD sparse matrix for a given transverse
periodic shift.  Along the transport direction $z$:
\begin{enumerate}
  \item The supercell size in the transport direction is $G_z / 2$
        (integer division), and 1 in the transverse directions.  For
        $G_z = 3$, this gives a supercell of one unit cell
        ($3 // 2 = 1$).
  \item Three BTD blocks are extracted for each transverse periodic
        shift $(n_x,n_y)$:
        \begin{itemize}
          \item $n_z=-1$: lower off-diagonal,
          \item $n_z=0$: diagonal,
          \item $n_z=+1$: upper off-diagonal.
        \end{itemize}
  \item Each block is assembled by \texttt{\_get\_transport\_block},
        which loops over all pairs of local shifts within the supercell
        and selects the appropriate unit-cell block for each pair.
  \item The three blocks are tiled $N_\text{transport}$ times along the
        diagonal to build the full BTD sparse matrix.
\end{enumerate}

The function \texttt{\_assemble\_kpoint} then performs the Fourier sum
over transverse periodic shifts.  For each $\vc{k}_\perp$ on a
Monkhorst--Pack mesh:
\begin{equation}\label{eq:assemble}
  D(\vc{k}_\perp)
  = \sum_{(n_x,n_y)} M(n_x,n_y)\;
    e^{2\pi i\,(k_x n_x + k_y n_y)},
\end{equation}
where $M(n_x,n_y)$ is the BTD sparse matrix for periodic shift
$(n_x,n_y)$, already expanded in the transport direction.  The
$k$-points on the Monkhorst--Pack mesh are
\begin{equation}
  k_i = \frac{n_i + 0.5}{N_i} - 0.5 + \text{shift}_i,
  \qquad n_i = 0,\ldots,N_i-1.
\end{equation}
The result is a stack of BTD matrices indexed by
$(\omega,k_x,k_y)$.

% ======================================================================
\section{Harmonic Transport Validation}
\label{sec:harmonic_validation}
% ======================================================================

This section summarizes the harmonic transport problem solver in
QuaTrEx, the validation checks used to verify the model, and the
reference calculations used as baselines for later
anharmonic extensions.

% ----------------------------------------------------------------------
\subsection{Ballistic phonon NEGF}
\label{sec:negf}
% ----------------------------------------------------------------------

In the ballistic limit, phonon transport is described by the harmonic
dynamical matrix together with open boundary conditions representing the
semi-infinite contacts.  For each transverse wavevector
$\vc{k}_\perp$, the problem reduces to a one-dimensional transport
problem along the transport direction, here taken to be $z$.  The
relevant objects are the transport-direction block matrices, the contact
surface Green's functions, and the retarded and Keldysh Green's
functions of the device.

\subsubsection{Transport-direction block decomposition}

For a fixed transverse wavevector
$\vc{k}_\perp = (k_x,k_y)$, the on-site and nearest-neighbour coupling
blocks along the transport direction are obtained by a partial Fourier
transform over the $q_z$ mesh.  Defining the harmonic real-space blocks
in Convention~B as $H(n_x,n_y,n_z)$, the transport-direction blocks are
\begin{align}
  H_{00}(\vc{k}_\perp)
    &= \frac{C}{N_{q_z}}
       \sum_{k=0}^{N_{q_z}-1}
       D^{(\mathrm{B})}\!\left(k_x, k_y, \frac{k}{N_{q_z}}\right),
  \label{eq:H00} \\
  H_{01}(\vc{k}_\perp)
    &= \frac{C}{N_{q_z}}
       \sum_{k=0}^{N_{q_z}-1}
       D^{(\mathrm{B})}\!\left(k_x, k_y, \frac{k}{N_{q_z}}\right)
       e^{-2\pi i k / N_{q_z}}.
  \label{eq:H01}
\end{align}
At non-zero transverse wavevector, these quantities are generally
complex because they arise from a two-dimensional Fourier sum over the
real-space transverse shifts,
\begin{equation}
  H_{0n_z}(\vc{k}_\perp)
  = \sum_{n_x,n_y}
    H(n_x,n_y,n_z)\,e^{2\pi i(k_x n_x + k_y n_y)}.
\end{equation}
Hermiticity implies
\begin{equation}
  H_{10}(\vc{k}_\perp) = H_{01}^\dagger(\vc{k}_\perp).
\end{equation}

These blocks define the effective one-dimensional transport problem for
each transverse wavevector and form the input to the contact and device
Green's-function calculation.

\subsubsection{Sancho--Rubio decimation}

The surface Green's function of a semi-infinite harmonic lead is
computed iteratively using the Sancho--Rubio decimation method.  Given
the bulk on-site and coupling blocks $H_{00}$ and $H_{01}$, the
algorithm is initialized as
\begin{equation}
  \epsilon = \epsilon_s = (\omega^2 + i\eta)I - H_{00},
  \qquad
  \alpha = -H_{10},
  \qquad
  \beta = -H_{01},
\end{equation}
and then iterates
\begin{align}
  \epsilon_s &\leftarrow \epsilon_s - \alpha\,\epsilon^{-1}\beta, \\
  \epsilon   &\leftarrow \epsilon - \alpha\,\epsilon^{-1}\beta
              - \beta\,\epsilon^{-1}\alpha, \\
  \alpha     &\leftarrow \alpha\,\epsilon^{-1}\alpha, \\
  \beta      &\leftarrow \beta\,\epsilon^{-1}\beta.
\end{align}
The iteration is terminated when
\begin{equation}
  \|\alpha\| + \|\beta\| < \mathrm{tol},
\end{equation}
at which point the surface Green's function is
\begin{equation}
  g_s = \epsilon_s^{-1}.
\end{equation}

For the left lead, the surface Green's function is computed as
\begin{equation}
  g_L = g_s(H_{00},H_{01}),
\end{equation}
whereas for the right lead the coupling direction is reversed,
\begin{equation}
  g_R = g_s(H_{00},H_{01}^\dagger).
\end{equation}

\subsubsection{Caroli formula}

Once the lead surface Green's functions have been computed, the retarded
contact self-energies are
\begin{align}
  \Sigma_L &= H_{10}\,g_L\,H_{01}, &
  \Sigma_R &= H_{01}\,g_R\,H_{10},
\end{align}
and the corresponding broadening matrices are
\begin{equation}
  \Gamma_{L/R} = i(\Sigma_{L/R} - \Sigma_{L/R}^\dagger).
\end{equation}
The retarded device Green's function is then
\begin{equation}
  G^R
  = \bigl[(\omega^2+i\eta)I - H_D - \Sigma_L - \Sigma_R\bigr]^{-1},
\end{equation}
with advanced Green's function
\begin{equation}
  G^A = (G^R)^\dagger.
\end{equation}

The ballistic transmission at fixed transverse wavevector is given by
the Caroli formula,
\begin{equation}\label{eq:caroli}
  \mathcal{T}(\omega^2,\vc{k}_\perp)
  = \mathrm{Tr}\bigl[\Gamma_L\,G^R\,\Gamma_R\,G^A\bigr].
\end{equation}
The transverse-wavevector-averaged transmission is then
\begin{equation}
  \bar{\mathcal{T}}(\omega^2)
  = \frac{1}{N_{k_\perp}}
    \sum_{\vc{k}_\perp}
    \mathcal{T}(\omega^2,\vc{k}_\perp).
\end{equation}

\subsubsection{QuaTrEx implementation}
\label{sec:phonon_solver}

In QuaTrEx, the harmonic problem is solved by the
\texttt{PhononSolver}.  The solver assembles the inverse retarded
Green's-function matrix in the form
\begin{equation}\label{eq:system}
  \bigl[G^R(\omega^2,\vc{k}_\perp)\bigr]^{-1}
  = (\omega^2+i\eta)\,\mathbb{I}
    - D(\vc{k}_\perp)
    - \Sigma^R(\omega^2,\vc{k}_\perp),
\end{equation}
where $D(\vc{k}_\perp)$ is the BTD harmonic dynamical matrix assembled
from the \texttt{dynamical\_matrix.mat} input, and
$\Sigma^R = \Sigma_L + \Sigma_R$ in the ballistic case.

The contact lesser and greater self-energies encode the Bose--Einstein
occupation of the leads,
\begin{align}
  \Sigma_L^< &= i\,\Gamma_L\,n_\mathrm{BE}(\omega,T_L), &
  \Sigma_L^> &= i\,\Gamma_L\,\bigl[n_\mathrm{BE}(\omega,T_L)+1\bigr],
\end{align}
and analogously for the right contact.  The Keldysh Green's functions
are then obtained from
\begin{align}
  G^< &= G^R\,\Sigma_\mathrm{total}^<\,(G^R)^\dagger, \\
  G^> &= G^R\,\Sigma_\mathrm{total}^>\,(G^R)^\dagger.
\end{align}

This harmonic implementation is the baseline solver used directly in the
ballistic calculations and also the zeroth-order problem in the later
SCBA treatment of anharmonic transport.

\subsubsection{Thermal conductance}

The ballistic thermal conductance per unit area is obtained from the
transmission through
\begin{equation}\label{eq:G_thermal}
  G(T)
  = \frac{1}{A_c}
    \int_0^\infty \frac{d\omega}{2\pi}\;
    \bar{\mathcal{T}}(\omega)\;
    \hbar\omega\;
    \frac{\partial n_\mathrm{BE}}{\partial T},
\end{equation}
where $A_c$ is the cross-sectional area perpendicular to transport and
$n_\mathrm{BE}(\omega,T)$ is the Bose--Einstein distribution.  Its
temperature derivative is
\begin{equation}
  \frac{\partial n_\mathrm{BE}}{\partial T}
  = \frac{\hbar\omega}{k_B T^2}\,
    \frac{e^{\hbar\omega/k_B T}}
         {(e^{\hbar\omega/k_B T}-1)^2}.
\end{equation}
In practice, the integral is evaluated numerically on the sampled
frequency grid using the trapezoidal rule.

% ----------------------------------------------------------------------
\subsection{Validation of the extracted harmonic model}
\label{sec:extraction_validation}
% ----------------------------------------------------------------------

Before using the extracted harmonic blocks in transport calculations, it
is important to verify that the block representation is internally
consistent and reproduces the original phonopy dynamical matrix.  The
present validation procedures are implemented in
\texttt{validation.py}.

\paragraph{Acoustic sum rule.}

At the $\Gamma$ point, the lattice Fourier series reduces to
\begin{equation}
  D^{(\mathrm{B})}(\vc{0}) = \sum_{\vc{R}} H(\vc{R}).
\end{equation}
Translational invariance requires that the three acoustic modes vanish
at $\Gamma$.  The validation therefore forms
\begin{equation}
  D_\Gamma = \sum_{\vc{R}} H(\vc{R}),
\end{equation}
computes its eigenvalues $\lambda_n$, and reports the corresponding
frequencies
\begin{equation}
  \omega_n = \mathrm{sign}(\lambda_n)\sqrt{|\lambda_n|}.
\end{equation}
The three smallest $|\omega_n|$ should be close to zero.  In addition,
the matrix should be symmetric to numerical precision, so that
\begin{equation}
  \max |D_\Gamma - D_\Gamma^T|
\end{equation}
is approximately zero.

\paragraph{Block symmetry.}

The real-space harmonic blocks must satisfy the Hermiticity relation
\begin{equation}\label{eq:block_sym}
  H(\vc{R})^T = H(-\vc{R})
  \qquad \forall\;\vc{R}.
\end{equation}
The validation iterates over all block pairs
$(\vc{R},-\vc{R})$ and reports the maximum deviation
\begin{equation}
  \max_{\vc{R}}
  \|H(\vc{R})^T - H(-\vc{R})\|_\infty.
\end{equation}

\paragraph{Band-structure comparison.}

A more stringent validation compares phonon frequencies reconstructed
from the extracted blocks with those obtained directly from phonopy.
Along a chosen $\vc{q}$-path, two spectra are computed:
\begin{enumerate}
  \item the phonopy reference spectrum, obtained from
        \begin{equation}
          D_\mathrm{ph}(\vc{q})
          = P(\vc{q})\,D^{(\mathrm{A})}(\vc{q})\,P^\dagger(\vc{q}),
        \end{equation}
        converted to $(\mathrm{rad/s})^2$ and diagonalized;

  \item the reconstructed spectrum, obtained from
        \begin{equation}
          D_\mathrm{rec}(\vc{q})
          = \sum_{\vc{R}} H(\vc{R})\,
            e^{2\pi i\,\vc{q}\cdot\vc{R}},
        \end{equation}
        and then diagonalized.
\end{enumerate}
The maximum eigenfrequency deviation
\begin{equation}
  \max_{n,\vc{q}}
  |\omega_n^\mathrm{ph} - \omega_n^\mathrm{rec}|
\end{equation}
serves as a direct measure of the quality of the extraction.  In
addition, overlay plots of the two band structures provide a visual
consistency check.

% ----------------------------------------------------------------------
\subsection{Heterogeneous-device extension}
\label{sec:heterogeneous}
% ----------------------------------------------------------------------

The original harmonic unit-cell input mode in QuaTrEx constructs a
homogeneous device by repeating the same block set across all transport
cells.  To model interfaces such as Si/Ge, the harmonic assembly was
extended to support piecewise-defined left, device, and right regions
with distinct block dictionaries.

\subsubsection{L|D|R device model}

The transport region is partitioned into three sections,
\begin{equation}\label{eq:ldr}
  \underbrace{L\;\cdots\;L}_{n_L\;\text{cells}}
  \;\underbrace{D\;\cdots\;D}_{n_D\;\text{cells}}
  \;\underbrace{R\;\cdots\;R}_{n_R\;\text{cells}},
  \qquad
  n_L + n_D + n_R = n_\mathrm{transport},
\end{equation}
where the left and right parts serve as lead-like material regions and
the central part represents the interface or device region.

\subsubsection{Configuration changes}

To support this heterogeneous harmonic setup, four new configuration
fields were added to \texttt{DeviceConfig}:
\begin{center}
\begin{tabular}{lll}
  \toprule
  Field & Type & Purpose \\
  \midrule
  \texttt{num\_left\_cells}  & \texttt{PositiveInt | None} & number of left cells $n_L$ \\
  \texttt{num\_right\_cells} & \texttt{PositiveInt | None} & number of right cells $n_R$ \\
  \texttt{left\_matrix}      & \texttt{str | None}         & left-lead \texttt{.mat} file \\
  \texttt{right\_matrix}     & \texttt{str | None}         & right-lead \texttt{.mat} file \\
  \bottomrule
\end{tabular}
\end{center}
This mode requires
\texttt{construct\_from\_unit\_cell = true} and must satisfy
\begin{equation}
  n_L + n_R < n_\mathrm{transport},
\end{equation}
so that at least one transport cell remains in the device region.

\subsubsection{Matrix assembly}

The heterogeneous assembly is implemented by extending the harmonic
unit-cell expansion to allow different block dictionaries in different
transport cells.  Two new functions were added to
\texttt{quatrex/device/inputs.py}.

\texttt{\_expand\_heterogeneous\_matrix} builds a BTD sparse matrix for
a single transverse periodic shift $(n_x,n_y)$.  Each cell is assigned
to one of the three regions according to its index.  The diagonal block
of cell $i$ is selected according to
\begin{equation}
  D_{ii} =
  \begin{cases}
    D_L, & i < n_L, \\
    D_D, & n_L \le i < n_L+n_D, \\
    D_R, & i \ge n_L+n_D.
  \end{cases}
\end{equation}

\texttt{\_create\_heterogeneous\_matrix} loops over all transverse
periodic shifts, calls \texttt{\_expand\_heterogeneous\_matrix} for
each, and combines them into the final transport object.  At $\Gamma$,
the transverse contributions are also summed to determine the sparsity
pattern.

\subsubsection{Boundary coupling}

At the interfaces between the $L$, $D$, and $R$ regions, the
off-diagonal couplings must be chosen consistently with Hermiticity.
For harmonic real-space blocks, the relevant relation is
\begin{equation}\label{eq:hermiticity_perp}
  H(\vc{R}_\perp,n_z)^\dagger
  = H(-\vc{R}_\perp,-n_z).
\end{equation}
At the left/device boundary, the super-diagonal coupling is taken from
the left-lead block $H_{01}^L$, while the corresponding sub-diagonal is
$H_{10}^L$.  At the device/right boundary, the couplings are taken from
the right-lead blocks $H_{01}^R$ and $H_{10}^R$.  This ensures that the
assembled heterogeneous matrix remains consistent with the physical lead
couplings on either side of the interface.

% ----------------------------------------------------------------------
\subsection{Reference calculations}
\label{sec:validation}
% ----------------------------------------------------------------------

To validate the harmonic transport implementation, the QuaTrEx results
are compared against simpler reference calculations based on direct
Sancho--Rubio decimation and the Caroli formula.  A naive
Sancho--Rubio + Caroli Python code is used as reference.  These
reference cases serve both as correctness checks for the harmonic
backbone and as a baseline for interpreting later anharmonic
corrections.

\subsubsection{Bulk validation}
\label{sec:bulk_validation}

The simplest test is a homogeneous bulk silicon system.  In this case,
the extracted harmonic blocks are used both in QuaTrEx and in a
standalone reference implementation of Sancho--Rubio plus Caroli.  Both
calculations are run with identical harmonic inputs, the same
transverse-wavevector mesh, and the same frequency grid.

Agreement between the two transmission curves confirms that:
\begin{itemize}
  \item the block extraction from phonopy is correct,
  \item the BTD assembly from unit-cell blocks is consistent,
  \item the open-boundary-condition treatment in QuaTrEx matches the
        reference implementation.
\end{itemize}

\subsubsection{Si/Ge interface}
\label{sec:sige}

\begin{figure}[htb]
  \centering
  \includegraphics[width=\textwidth]{fig/reference/si_ge_quatrex_comparison.png}
  \caption{Ballistic phonon transmission on a $20\times20$ transverse
    $k$-mesh with 101 frequency points.
    \textbf{Left:} Homogeneous bulk Si.
    \textbf{Right:} Si/Ge interface in the mass-mismatch approximation.}
  \label{fig:quatrex_comparison}
\end{figure}

As a more realistic reference problem, a ballistic Si/Ge interface is
constructed using a mass-mismatch model.  In this approximation, the
harmonic force constants are taken from bulk Si throughout, and the
distinction between Si and Ge enters only through the atomic masses,
\begin{equation}
  m_\mathrm{Si} = 28.09\;\mathrm{amu},
  \qquad
  m_\mathrm{Ge} = 72.64\;\mathrm{amu}.
\end{equation}
This is justified by the similarity of the Si and Ge interatomic force
constants together with their substantial mass difference.

The heterogeneous system is partitioned into three regions:
\begin{itemize}
  \item \textbf{Left lead:} Si force constants with Si masses,
  \item \textbf{Device:} Si force constants with mixed Si/Ge masses,
  \item \textbf{Right lead:} Si force constants with Ge masses.
\end{itemize}
In the present setup, atoms with fractional coordinate $z<0.5$ are
assigned Si mass in the device cell and the remainder Ge mass.  The
coupling between left lead and device is taken from the Si lead, while
the coupling between device and right lead is taken from the Ge lead.

For the harmonic inputs, the Si force constants are obtained from a DFT
calculation at equilibrium lattice constant
\begin{equation}
  a = 5.466\;\mathrm{\AA},
\end{equation}
using the conventional 8-atom diamond cell and a
$2\times2\times2$ supercell.  The DFT parameters are:
plane-wave cutoff 60\,Ry, charge-density cutoff 480\,Ry,
$4\times4\times4$ supercell $k$-mesh, SCF threshold
$10^{-10}$\,Ry, Gaussian smearing 0.01\,Ry, PBE exchange-correlation,
and the pseudopotential
\texttt{Si.pbe-n-rrkjus\_psl.1.0.0.UPF}.  Because of diamond symmetry,
this setup reduces to a single displaced calculation in the harmonic
finite-displacement workflow.

The corresponding phonopy calculation is performed with the native QE
interface, so the lattice and force constants are converted from bohr
and Ry/bohr$^2$ to \AA\ and eV/\AA$^2$ on loading, as described in
\cref{sec:phonopy_units}.  The Ge and interface models are obtained by
cloning the Si harmonic model and substituting the atomic masses.

The transport parameters used for the Si/Ge interface calculation are
\begin{center}
\begin{tabular}{ll}
  \toprule
  Parameter & Value \\
  \midrule
  Transverse $k$-mesh & $20\times20$ \\
  Frequency range & $0.01$--$16.0$\,THz, 101 points \\
  Transport direction & $z$ (001) \\
  Total transport cells / $n_L$ / $n_D$ / $n_R$ & 5 / 2 / 1 / 2 \\
  Broadening $\eta$ & $\tfrac12 \Delta\omega^2 \approx 5.05\times10^{23}\,(\mathrm{rad/s})^2$ \\
  $T_L$, $T_R$ & 301\,K, 299\,K \\
  OBC algorithm & Sancho--Rubio \\
  \bottomrule
\end{tabular}
\end{center}

At the $\Gamma$ point, the maximum optical frequency of the Si model is
15.55\,THz, consistent with reference values.  The Ge mass-mismatch
clone gives 9.67\,THz, in agreement with the simple mass-scaling
relation
\begin{equation}
  f_\mathrm{Ge}
  = f_\mathrm{Si}\sqrt{\frac{m_\mathrm{Si}}{m_\mathrm{Ge}}}.
\end{equation}

The resulting interfacial thermal conductance at 300\,K is
\begin{equation}
  G = 196\;\mathrm{MW/(m^2\,K)}.
\end{equation}
This may be compared with the literature values
\begin{center}
\begin{tabular}{ll}
  \toprule
  Reference & $G$ [MW/(m$^2$\,K)] \\
  \midrule
  Latour et al.\ (2017) & 210 \\
  Tian et al.\ (2012) & 209 \\
  Present calculation & 196 \\
  \bottomrule
\end{tabular}
\end{center}
The remaining $\sim 7\%$ deviation is likely attributable to the coarser
transverse-wavevector mesh and the specific DFT setup used here.

% ======================================================================
\section{Anharmonic Extension: FC3 and Self-Consistent Transport}
\label{sec:anharmonic_transport}
% ======================================================================

The harmonic transport problem described in the previous section treats
phonons as non-interacting quasiparticles propagating through the device
under the action of the quadratic force-constant matrix alone.  The
anharmonic extension adds phonon--phonon scattering through the cubic
interatomic force constants (FC3), represented here by the mass-weighted
tensor $\Psi$.  In practice, this turns the ballistic transport problem
into a self-consistent many-body problem: the Green's functions depend
on the phonon self-energy, while the phonon self-energy itself is
constructed from the Green's functions.

The present implementation proceeds in three layers.  First, the full
FC3 dataset is reduced to a transport-compatible local tensor.  Second,
this tensor is used to construct the phonon--phonon self-energy at
lowest nontrivial order in the cubic interaction.  Third, the resulting
retarded, lesser, and greater self-energies are incorporated into a
self-consistent Born approximation (SCBA) loop.  The remainder of this
section summarizes the reduced FC3 models currently used, the basic
self-energy formulas, the SCBA procedure, and more advanced
formulations and scaling analyses relevant for large-scale distributed
calculations.

% ----------------------------------------------------------------------
\subsection{Reduced FC3 models for transport}
\label{sec:assembly_fc3}
% ----------------------------------------------------------------------

In the present transport implementation, the harmonic device region is
one transport cell thick in the transport direction, with harmonic
dimension
\begin{equation}
  n_\mathrm{dof} = 3N_b,
\end{equation}
where $N_b$ is the number of basis atoms in the transport cell.  For
the conventional 8-atom silicon cell used here, $n_\mathrm{dof}=24$.
The full real-space FC3 dataset, however, couples triplets of atoms
across different unit cells, and therefore exceeds the locality of the
harmonic transport slab representation.

To obtain a tractable transport model, the full FC3 dataset is reduced
to an effective local tensor $\Psi_{abc}$ for the transport slab.  Two
approximations are considered below.

\subsubsection{Onsite approximation}

The simplest reduction retains only the fully local FC3 block, i.e. the
contribution for which both relative cell vectors are zero:
\begin{equation}
  \Psi_{abc}^{\mathrm{onsite}}
  =
  \frac{\Phi^{(3)}_{abc}(\vc{0},\vc{0},\vc{0})}
       {\sqrt{M_a M_b M_c}}.
\end{equation}
Here $a,b,c$ denote the combined Cartesian and basis indices of the
transport slab.  This approximation ignores all cubic couplings to
neighbouring transport slabs and retains only scattering processes that
are fully local within the reference slab.

\subsubsection{Full \texorpdfstring{$\delta l = 0$}{delta l = 0} approximation}

A less restrictive approximation retains all FC3 blocks whose two
relative cell vectors remain within the same transport slab as the
reference atom.  This corresponds to the condition
\begin{equation}
  \delta l' = \delta l'' = 0,
\end{equation}
where the slab displacement along the transport direction is extracted
from the real-space cell vectors through
\begin{equation}
  \delta l'
  = \mathrm{round}\!\bigl([\vc{A}^{-\top}\vc{R}']_z\bigr),
\end{equation}
and similarly for $\delta l''$.  The reduced slab tensor is then
\begin{equation}\label{eq:fc3_full}
  \Psi_{abc}^{\mathrm{full}}
  =
  \sum_{\substack{\vc{R}',\vc{R}'' \\
                  \delta l' = 0,\;\delta l'' = 0}}
  \frac{\Phi^{(3)}_{abc}(\vc{0},\vc{R}',\vc{R}'')}
       {\sqrt{M_a M_b M_c}}.
\end{equation}
This is the transport-slab analogue of Approximation~III in the
literature on anharmonic phonon transport.  It retains all cubic
interactions whose two partner atoms remain in the same transport slab,
while still discarding couplings to neighbouring slabs.

For the silicon FC3 dataset used here, 193 of the 422 remapped FC3
blocks satisfy the condition $\delta l'=\delta l''=0$, corresponding to
2126 nonzero tensor elements out of the total
$24^3 = 13\,824$ possible entries.  This provides a significantly richer
local anharmonic model than the onsite approximation while remaining
compatible with the current single-slab transport treatment.

% ----------------------------------------------------------------------
\subsection{Phonon--phonon self-energy}
\label{sec:self_energy}
% ----------------------------------------------------------------------

Given the reduced cubic interaction tensor $\Psi$, the lowest-order
phonon--phonon scattering self-energy is constructed from a pair of
Green's functions connected by two FC3 vertices.  This is the phonon
analogue of a second-order Born self-energy.  In the present context,
the lesser and greater components are the primary quantities, and the
retarded self-energy is approximated from them.

\subsubsection{Local self-energy formula}

For a local transport slab, the phonon--phonon self-energy at second
order in the cubic interaction is
\begin{equation}\label{eq:sigma_guo}
  \Sigma^{<}_{ab}(\omega)
  = \frac{i\hbar}{2}
    \sum_{c,d,e,f}
    \Psi_{acd}
    \left[
      \int \frac{d\omega'}{2\pi}\;
      G^<_{ce}(\omega')\,
      G^<_{df}(\omega-\omega')
    \right]
    \Psi_{bef},
\end{equation}
where $G^<(\omega)$ is the slab lesser Green's function and
$\Psi$ is the mass-weighted FC3 tensor in SI units.  The same structure
applies to the greater self-energy,
\begin{equation}
  \Sigma^{>}_{ab}(\omega)
  = \frac{i\hbar}{2}
    \sum_{c,d,e,f}
    \Psi_{acd}
    \left[
      \int \frac{d\omega'}{2\pi}\;
      G^>_{ce}(\omega')\,
      G^>_{df}(\omega-\omega')
    \right]
    \Psi_{bef}.
\end{equation}

In the current implementation, the retarded phonon self-energy is
approximated by neglecting the Hilbert-transform contribution and using
\begin{equation}\label{eq:sigma_r}
  \Sigma^R(\omega)
  \approx \frac{1}{2}
  \bigl[\Sigma^>(\omega) - \Sigma^<(\omega)\bigr].
\end{equation}
This approximation captures the imaginary, scattering-related part of
the retarded self-energy, but it does not yet include the full
dispersive real-part correction.

\subsubsection{Full Fourier-space self-energy}

For a genuinely three-dimensional system with periodicity transverse to
the transport direction, the phonon--phonon self-energy depends on the
transverse wavevector and includes an internal sum over all transverse
momenta.  The full lesser/greater self-energy takes the form
\begin{align}\label{eq:sigma_phph}
  \Sigma^{\gtrless}_{s,\,ll'_s}(&\omega;\,\vc{q}_\perp)
  = \frac{i\hbar}{2}
    \sum_{\substack{l_1 l_2 l_3 l_4\\ j_1 j_2 j_3 j_4}}
    \frac{1}{N}\sum_{\vc{q}'_\perp}
    \tilde{\Phi}^{ij_1j_2}_{l,\,l_1 l_2}
      \!\left(\vc{q}'_\perp,\vc{q}_\perp-\vc{q}'_\perp\right)
    \notag\\
  &\qquad\times
    \tilde{\Phi}^{j_3j_4j}_{l',\,l_3 l_4}
      \!\left(\vc{q}'_\perp-\vc{q}_\perp,-\vc{q}'_\perp\right)
    \int\!\frac{d\omega'}{2\pi}\;
    G^{\gtrless,j_1j_4}_{l_1l_4}(\omega';\vc{q}'_\perp)\,
    G^{\gtrless,j_2j_3}_{l_2l_3}
      (\omega-\omega';\vc{q}_\perp-\vc{q}'_\perp).
\end{align}
Here $l,l'$ index transport slabs, $s$ indicates local basis indices
within a slab, and the sum over $\vc{q}'_\perp$ enforces transverse
momentum conservation.

The Fourier-transformed FC3 tensor entering this expression is
\begin{equation}\label{eq:fc3_fourier}
  \tilde{\Phi}^{ij_1j_2}_{l,\,l_1 l_2}
  (\vc{q}_\perp,\vc{q}'_\perp)
  = \sum_{\Delta\vc{R}_\perp}\sum_{\Delta\vc{R}'_\perp}
    \Phi^{ij_1j_2}_{l l_1 l_2}\;
    e^{-i\vc{q}_\perp\cdot\Delta\vc{R}_\perp}\;
    e^{-i\vc{q}'_\perp\cdot\Delta\vc{R}'_\perp}.
\end{equation}
Unlike the harmonic FC2 object, which depends on one transverse
wavevector, the cubic tensor depends on two.  This is the origin of the
substantially greater computational complexity of the anharmonic
problem.

In the ballistic harmonic case, each $\vc{q}_\perp$ is independent and
the transport problem is embarrassingly parallel over wavevectors.  By
contrast, the internal $\vc{q}'_\perp$ sum in
\cref{eq:sigma_phph} couples all transverse wavevectors, so that
evaluating the self-energy at one $\vc{q}_\perp$ requires access to
Green's functions at every other transverse wavevector.

\subsubsection{Frequency discretisation and FFT convolution}
\label{sec:fft}

In practice, the self-energy is evaluated on a discrete real-frequency
grid
\begin{equation}
  \omega_n = \omega_\mathrm{min} + n\,\Delta\omega,
  \qquad
  n = 0,\ldots,N_\omega-1.
\end{equation}
The broadening parameter is chosen as
\begin{equation}
  \eta = \eta_\mathrm{factor}\,(\Delta\omega)^2,
\end{equation}
with $\eta_\mathrm{factor}=0.5$ in the current implementation.

The frequency integral appearing in \cref{eq:sigma_guo} is a
convolution.  On the discrete frequency grid, it is approximated by
\begin{equation}
  C_{ce,df}[\omega_n]
  = \frac{\Delta\omega}{2\pi}
    \sum_m
    G^<_{ce}[\omega_m]\;
    G^<_{df}[\omega_n-\omega_m].
\end{equation}
This convolution is evaluated efficiently by FFT:
\begin{equation}
  C
  = \mathrm{IFFT}\!\bigl[
      \mathrm{FFT}(G_{ce})\cdot \mathrm{FFT}(G_{df})
    \bigr].
\end{equation}
The same procedure is applied to the greater component.

The resulting discrete lesser self-energy is
\begin{equation}\label{eq:full-energy}
  \Sigma^<_{ab}[\omega_n]
  = \frac{i\hbar\,\Delta\omega}{4\pi}
    \sum_{c,d,e,f}
    \Psi_{acd}\;C_{ce,df}[\omega_n]\;\Psi_{bef},
\end{equation}
and analogously for $\Sigma^>$.  In the current code this is
implemented as a loop over active $(c,d)$ and $(e,f)$ index pairs with
sparsity checks on the FC3 tensor.  For the silicon test case, the
number of active pairs is of order $10^4$, leading to an overall
self-energy cost scaling approximately as
\begin{equation}
  \mathcal{O}(n_\mathrm{dof}^4 N_\omega \log N_\omega).
\end{equation}

% ----------------------------------------------------------------------
\subsection{SCBA algorithm}
\label{sec:scba}
% ----------------------------------------------------------------------

The phonon self-energy modifies the retarded and Keldysh Green's
functions, which in turn modify the self-energy itself.  The resulting
nonlinear problem is solved iteratively within the self-consistent Born
approximation (SCBA).  Starting from the harmonic ballistic solution,
the phonon self-energy is updated repeatedly until convergence.

\subsubsection{SCBA loop}

The current SCBA procedure is:
\begin{enumerate}
  \item \textbf{Initialization.}
        Set the phonon self-energy to zero,
        \begin{equation}
          \Sigma^\mathrm{ph} = 0,
        \end{equation}
        and solve the harmonic transport problem.

  \item \textbf{Green's functions.}
        For each transverse wavevector $\vc{q}_\perp$ and frequency
        $\omega_n$, compute the contact self-energies using
        Sancho--Rubio decimation and solve
        \begin{align}
          G^R
          &= \Bigl[
               (\omega^2+i\eta)\,\mathbb{I}
               - H_{00}
               - \Sigma_L^R
               - \Sigma_R^R
               - \Sigma_\mathrm{ph}^R
             \Bigr]^{-1}, \\
          G^{<,>}
          &= G^R\,
             \Sigma^{<,>}_\mathrm{total}\,
             (G^R)^\dagger.
        \end{align}

  \item \textbf{Self-energy update.}
        Construct the new
        $\Sigma_\mathrm{ph}^{<,>,R}$ from the Green's functions using
        the discretized self-energy formulas of
        \cref{eq:full-energy,eq:sigma_r}.

  \item \textbf{Mixing.}
        Stabilize the iteration by linear mixing,
        \begin{equation}
          \Sigma^{(n+1)}
          = (1-\alpha)\Sigma^{(n)}
            + \alpha\,\Sigma^\mathrm{new},
        \end{equation}
        with default mixing parameter $\alpha=0.3$.

  \item \textbf{Convergence check.}
        Evaluate a convergence criterion based on the change in the
        transport observable between two iterations.  In the current
        implementation this is written as
        \begin{equation}
          \frac{\sum_n
            |\mathcal{T}^{(k)}_n-\mathcal{T}^{(k-1)}_n|}
               {\sum_n |\mathcal{T}^{(k-1)}_n|}
          < \mathrm{tol},
        \end{equation}
        although the precise transport quantity used here is still under
        review.

  \item \textbf{Iteration.}
        Repeat until the solution converges or the maximum number of
        iterations is reached.
\end{enumerate}

Per SCBA iteration, the dominant costs are:
\begin{itemize}
  \item the solution of $N_\perp\times N_\omega$ Green's-function
        problems, including Sancho--Rubio evaluations for the contacts,
  \item the FFT-based self-energy update,
        \begin{equation}
          \mathcal{O}(n_\mathrm{dof}^4 N_\omega\log N_\omega),
        \end{equation}
  \item the evaluation of the transport observable used in the
        convergence criterion.
\end{itemize}
For sufficiently large transport slabs, the self-energy update is
expected to dominate the iteration cost.

\subsubsection{Implementation status and current limitations}

The present SCBA implementation should be regarded as a working
prototype rather than a final production method.  Several aspects are
already operational, while others remain approximate or still require
revision.

The following components are implemented:
\begin{itemize}
  \item FC3 remapping to the conventional harmonic basis,
  \item local reduced FC3 models for the transport slab,
  \item FFT-based evaluation of the lesser and greater self-energy,
  \item inclusion of the phonon self-energy into the retarded and
        Keldysh Green's-function equations,
  \item iterative SCBA updates with linear mixing.
\end{itemize}

The following limitations remain:
\begin{itemize}
  \item The retarded phonon self-energy currently uses the approximate
        relation \cref{eq:sigma_r} and does not yet include the full
        Hilbert-transform contribution.
  \item The transport quantity used inside the SCBA convergence loop is
        not yet finalized, and the corresponding transmission update
        logic still needs revision.
  \item The current transport implementation retains only slab-local FC3
        contributions, i.e. onsite or $\delta l=0$ approximations.
        Fully nonlocal anharmonic couplings along the transport
        direction are not yet included.
  \item The current implementation is mainly validated on bulk-like
        silicon test systems and has not yet been benchmarked as
        extensively as the harmonic solver.
\end{itemize}

Accordingly, the present anharmonic results should be interpreted as
preliminary and mainly as evidence that the overall SCBA machinery is
working and producing physically plausible trends.

\subsubsection{Preliminary anharmonic results}
\label{sec:anharmonic_results}

The current test system is bulk silicon in the conventional 8-atom
cell, with lattice constant
\begin{equation}
  a = 5.4662\;\mathrm{\AA}.
\end{equation}
The harmonic FC2 are obtained from a phonopy/QE
$2\times2\times2$ supercell calculation, while FC3 are taken from
\texttt{thirdorder.py} in the 2-atom FCC cell and remapped to the
conventional basis.  Transport is along the $z$ direction at 300\,K,
using 51 frequency points in the range 0.5--15.0\,THz, a mixing
parameter $\alpha=0.3$, and the slab-local FC3 approximation described
above.

The current results for two transverse $q$-meshes are summarized in
\cref{tab:scaling_anh}.
\begin{table}[h]
\centering
\caption{Self-energy magnitude and thermal conductance as a function of
the transverse $q$-mesh for the current silicon SCBA test.}
\label{tab:scaling_anh}
\begin{tabular}{cccccc}
  \toprule
  $q$-mesh & $N_\perp$ &
  $|\Sigma^R|_\mathrm{max}/|H_{00}|_\mathrm{max}$ &
  SCBA iter &
  $G_\mathrm{ball}$ &
  $G_\mathrm{anh}$ \\
  & & & & [MW/(m$^2$K)] & [MW/(m$^2$K)] \\
  \midrule
  $2\times2$ & 4  & 0.49 & 11 & 1096 & 913 \\
  $4\times4$ & 16 & 0.20 & 8  & 1162 & 1005 \\
  \bottomrule
\end{tabular}
\end{table}

These results show the expected suppression of thermal conductance in
the anharmonic calculation relative to the ballistic one.  The ratio
\begin{equation}
  \frac{|\Sigma^R|_\mathrm{max}}{|H_{00}|_\mathrm{max}}
\end{equation}
decreases as the transverse mesh is refined, indicating that the
effective self-energy per sampled point becomes less dominant once the
transverse phase space is better resolved.

The current transmission curves indicate that acoustic modes below
roughly 6\,THz are comparatively weakly affected, with transmission
ratio close to unity, whereas optical modes in the 10--15\,THz range
show a stronger reduction.  This trend is physically plausible: the
optical sector has larger phase space for three-phonon processes, while
the acoustic modes remain more constrained by energy conservation.
However, since the transmission post-processing is still under revision,
these frequency-resolved results should at present be interpreted only
qualitatively.

\begin{figure}[h]
  \centering
  \includegraphics[width=\textwidth]{fig/reference/anharmonic_test.png}
  \caption{Left: ballistic (blue) and anharmonic SCBA (red dashed)
    phonon transmission for bulk Si at 300\,K ($4\times4$ $q$-mesh,
    51 frequency points).
    Right: current transmission ratio
    $\mathcal{T}_\mathrm{anh}/\mathcal{T}_\mathrm{ball}$.}
  \label{fig:transmission}
\end{figure}

% ----------------------------------------------------------------------
\subsection{Advanced formulations of the phonon--phonon self-energy}
\label{sec:phph-energy}
% ----------------------------------------------------------------------

% ============================================================
\subsection{Notation and conventions}
\label{sub:notation}
% ============================================================

Throughout this section we use the following indices.
\begin{center}
\begin{tabular}{llr}
  \toprule
  Symbol & Meaning & Range \\
  \midrule
  $l,l'$               & Transport-direction slab index              & $N_\mathrm{sl}$ \\
  $s,s'$               & Basis-atom within a slab cross-section      & $N_s$~per~slab \\
  $i,j,j_1,\ldots,j_4$ & Cartesian component ($x,y,z$)               & $C=3$ \\
  $\vq_\perp$          & Transverse wavevector (2D BZ)                & $N_q$ points \\
  $\vR_\perp$          & Transverse lattice vector                    & $N_q$ points (dual) \\
  $\omega,\omega'$     & Real frequency                               & $N_\omega$ points \\
  $\tau$               & Conjugate time variable                      & $N_\tau = N_\omega$ \\
  $N$                  & Total number of $\vq$-points, $N_q$         & \\
  \bottomrule
\end{tabular}
\end{center}

\noindent
We define the combined Cartesian--basis index $\mu \equiv (j,s)$
with dimension $M \equiv CN_s = 3N_s$ per slab, and denote by $N_B$
the slab-pair bandwidth of the Green's function
($G^{\mu\mu'}_{ll'} = 0$ for $|l-l'|>N_B$).  The FC3 range in the
transport direction is $\hat{N}_c \equiv 2N_c+1$ slabs.

For each slab $l$, the transverse lattice vector is
$\vR_\perp(l) \equiv l_y\,\va_2 + l_z\,\va_3$,
where $(l_y, l_z)\in\mathbb{Z}^2$ and $\va_2, \va_3$ are the
transverse primitive vectors.  We define the transverse displacement
vectors
\begin{align}
  \vrho_1 &\equiv \vR_\perp(l)  - \vR_\perp(l_1), \label{eq:rho1}\\
  \vrho_2 &\equiv \vR_\perp(l)  - \vR_\perp(l_2), \label{eq:rho2}\\
  \vrho_3 &\equiv \vR_\perp(l') - \vR_\perp(l_3), \label{eq:rho3}\\
  \vrho_4 &\equiv \vR_\perp(l') - \vR_\perp(l_4). \label{eq:rho4}
\end{align}
With an FC3 cutoff of $N_\mathrm{cut}$ unit cells, each
displacement satisfies $|\vrho_k|\le N_\mathrm{cut}|\va|$, giving
at most $N_\perp \equiv (2N_\mathrm{cut}+1)^2$ nonzero values per
pair.

We define discrete Fourier transforms on the
$N_q = N_{q_2}\times N_{q_3}$ transverse periodic supercell:
\begin{align}
  G(\omega;\vq) &= \sum_{\vR} g(\omega;\vR)\,e^{-i\vq\cdot\vR},
  \label{eq:dft_fwd}\\
  g(\omega;\vR) &= \frac{1}{N}\sum_{\vq} G(\omega;\vq)\,
                    e^{+i\vq\cdot\vR}.
  \label{eq:dft_inv}
\end{align}
The discrete convolution theorem states
\begin{equation}
  \frac{1}{N}\sum_{\vq'} A(\vq')\,B(\vq-\vq')
  = \sum_{\vR} a(\vR)\,b(\vR)\,e^{-i\vq\cdot\vR},
  \label{eq:conv_theorem}
\end{equation}
where $a, b$ are the inverse DFTs of $A, B$.

The third-order force-constant (FC3)
$\Phi^{ij_1j_2}_{l,l_1,l_2}$ couples three atoms: one at slab $l$
with Cartesian index $i$, one at $(l_1, j_1)$, and one at
$(l_2, j_2)$.  Its transverse Fourier transform depends on two
wavevectors:
\begin{equation}
  \PhiT^{ij_1j_2}_{l,l_{1}l_{2}}(\vq,\vq')
  \;\equiv\;
  \sum_{\vrho_1}\sum_{\vrho_2}
  \Phi^{ij_1j_2}_{l,l_1,l_2}\;
  e^{-i\vq\cdot\vrho_1}\;e^{-i\vq'\cdot\vrho_2},
  \label{eq:fc3_fourier}
\end{equation}
where the sums run over all transverse displacements within the
cutoff and $\Phi^{ij_1j_2}_{l,l_1,l_2}$ is evaluated at the atom
positions selected by $\vrho_1$ and $\vrho_2$.


% ============================================================
\subsection{Self-Energy Transformations}
\label{sub:preliminary}
% ============================================================

The phonon--phonon self-energy is given by \cref{eq:sigma_phph},
\begin{align}
  \Sig^{\gtrless}_{s,ll'}(\omega;\vq)
  &= \frac{i\hbar}{2}
    \sum_{\substack{l_1 l_2 l_3 l_4\\j_1 j_2 j_3 j_4}}
    \frac{1}{N}\sum_{\vq'}
    \PhiT^{ij_1 j_2}_{l,\,l_1 l_2}\!\left(\vq',\vq{-}\vq'\right)
    \PhiT^{j_3 j_4 j}_{l',\,l_3 l_4}\!\left(\vq'{-}\vq,-\vq'\right)
    \notag\\
  &\qquad\times
    \int\!\dw\;
    G^{\gtrless,j_1j_4}_{l_1 l_4}(\omega';\vq')\,
    G^{\gtrless,j_2j_3}_{l_2 l_3}(\omega{-}\omega';\vq{-}\vq').
\end{align}
We suppress the $\gtrless$ superscript throughout.

% -------------------------------------------------------
\subsubsection{Time-domain representation of the frequency
  convolution}
\label{ssub:time_domain}
% -------------------------------------------------------

The $\omega'$-integral in \cref{eq:sigma_phph} is a convolution.
Define the fourier transforms
\begin{align}
  g(\tau;\vR) &= \frac{1}{N}\sum_{\vq}\int\!\frac{d\omega}{2\pi}\;
                  G(\omega;\vq)\,e^{-i\omega\tau}\,e^{+i\vq\cdot\vR},
  \label{eq:g_2d}\\
  G(\omega;\vq) &= \sum_{\vR}\int\!d\tau\;
                  g(\tau;\vR)\,e^{+i\omega\tau}\,e^{-i\vq\cdot\vR}.
  \label{eq:G_2d}
\end{align}
At fixed transverse argument the 1D convolution theorem gives
\begin{equation}
  \int\!\dw\; G_1(\omega';\cdot)\,G_2(\omega{-}\omega';\cdot)
  = \int\!d\tau\; g_1(\tau;\cdot)\,g_2(\tau;\cdot)\,
    e^{+i\omega\tau},
  \label{eq:omega_conv}
\end{equation}
replacing the $\omega'$-integration by a pointwise $\tau$-product
at cost $\cO(N_\omega\log N_\omega)$ per transverse point.

% -------------------------------------------------------
\subsubsection{Inner contraction}
\label{ssub:contraction}
% -------------------------------------------------------

At each working point, the self-energy accumulates a contribution
\begin{equation}
  \mathcal{B}^{ij}_{ll'}
  \;\mathrel{+}=\;
  \sum_{\substack{\mu_1 \mu_2 \mu_3 \mu_4\\l_1 l_2 l_3 l_4}}
  \Phi^{i\mu_1\mu_2}_{l,l_1,l_2}\;
  g^{\mu_1\mu_4}_{l_1,l_4}\;
  g^{\mu_2\mu_3}_{l_2,l_3}\;
  \Phi^{\mu_3\mu_4 j}_{l',l_3,l_4},
  \label{eq:inner_contraction}
\end{equation}
where $\Phi$ stands for either $\PhiT$ or the real-space FC3.
The four tensors form a ring: $\Phi_1$ shares $(\mu_1,l_1)$ with
$g_1$, which shares $(\mu_4,l_4)$ with $\Phi_2$, which shares
$(\mu_3,l_3)$ with $g_2$, which shares $(\mu_2,l_2)$ with
$\Phi_1$.  The optimal evaluation splits the ring into two halves:
\begin{align}
  {T_1}^{i\mu_2\mu_4}_{l,l_2,l_4}
  &= \sum_{\mu_1,l_1}
     \Phi^{i\mu_1\mu_2}_{l,l_1,l_2}\;g^{\mu_1\mu_4}_{l_1,l_4},
  \qquad
  {T_2}^{j\mu_2\mu_4}_{l',l_2,l_4}
  = \sum_{\mu_3,l_3}
     g^{\mu_2\mu_3}_{l_2,l_3}\;
     \Phi^{\mu_3\mu_4 j}_{l',l_3,l_4},
  \notag\\[4pt]
  \mathcal{B}^{ij}_{l,l'}
  &= \sum_{\mu_2,\mu_4,l_2,l_4} T_1\,T_2.
  \label{eq:balanced_tree}
\end{align}
Each half-transform contracts over $(\mu_1,l_1)$ of dimension
$M\hat{N}_c$, producing a result of dimension
$M\hat{N}_c \times MN_B$.  The cost is
\begin{equation}
  \cC_\mathrm{ic}
  = \cO\!\left(M^3\,\hat{N}_c^2\,N_B\right).
  \label{eq:cost_ic}
\end{equation}


% ============================================================
\subsection{\texorpdfstring{$\vq$}{q}-space formulation}
\label{sub:q_space}
% ============================================================

The $\vq$-space formulation retains the explicit $\vq'$-loop in
\cref{eq:sigma_phph} and converts the frequency convolution to a
$\tau$-product via \cref{eq:omega_conv}.

The Fourier-transformed FC3 is not stored but computed
on-the-fly at each $(\vq',\vq)$: one streams through the
$N_\perp^2$ nonzero real-space FC3 elements, multiplying each by
$e^{-i\vq'\cdot\vrho_1}\,e^{-i(\vq-\vq')\cdot\vrho_2}$ and
accumulating.  The phase computation costs
$\cO(N_\perp^2\,M^2\hat{N}_c^2)$ per $(\vq',\vq)$, which is
negligible compared to the ring contraction since
$N_\perp^2 \ll MN_B$.

Defining the $\tau$-dependent bubble at fixed $\vq$,
\begin{equation}
  \mathcal{B}(\tau;\vq)
  \;\equiv\;
  \frac{1}{N}\sum_{\vq'}
  \PhiT_1(\vq',\vq{-}\vq')\;
  g_1(\tau;\vq')\;
  g_2(\tau;\vq{-}\vq')\;
  \PhiT_2(\vq'{-}\vq,-\vq'),
  \label{eq:bubble_q}
\end{equation}
the self-energy is
\begin{equation}
  \Sig(\omega;\vq)
  = \frac{i\hbar}{2}
    \int\!d\tau\;
    \mathcal{B}(\tau;\vq)\,e^{+i\omega\tau}.
  \label{eq:sigma_q}
\end{equation}
At each $(\tau,\vq)$, the bubble is accumulated over $N_q$ values
of $\vq'$, each requiring one inner contraction
\cref{eq:balanced_tree}.  The total cost is
\begin{equation}
  \cW_q = N_\tau\,N_q^2\;\cC_\mathrm{ic}.
  \label{eq:cost_q}
\end{equation}


% ============================================================
\subsection{Real-space formulation}
\label{sub:real_space}
% ============================================================

We show that eliminating the $N_q^2$ loop by applying the
convolution theorem to the $\vq'$-sum is not competitive at the transverse
grids $N_q \approx 400$.

Expanding both Fourier-transformed FC3 factors in
\cref{eq:sigma_phph} using \cref{eq:fc3_fourier}:
\begin{align}
  \PhiT^{ij_1j_2}_{l,l_1l_2}\!\left(\vq',\vq-\vq'\right)
  &= \sum_{\vrho_1,\vrho_2}
     \Phi^{ij_1j_2}_{l,l_1,l_2}\;
     e^{-i\vq'\cdot\vrho_1}\;
     e^{-i(\vq-\vq')\cdot\vrho_2},
  \label{eq:v1_expand}\\[2pt]
  \PhiT^{j_3j_4j}_{l',l_3l_4}\!\left(\vq'-\vq,-\vq'\right)
  &= \sum_{\vrho_3,\vrho_4}
     \Phi^{j_3j_4j}_{l',l_3,l_4}\;
     e^{-i(\vq'-\vq)\cdot\vrho_3}\;
     e^{+i\vq'\cdot\vrho_4}.
  \label{eq:v2_expand}
\end{align}
The four exponentials factorize:
\begin{align}
  e^{-i\vq'\cdot\vrho_1}\, e^{-i(\vq-\vq')\cdot\vrho_2}
  &= e^{-i\vq\cdot\vrho_2}\cdot
     e^{-i\vq'\cdot(\vrho_1-\vrho_2)},
  \label{eq:v1_phases}\\[2pt]
  e^{-i(\vq'-\vq)\cdot\vrho_3}\, e^{+i\vq'\cdot\vrho_4}
  &= e^{+i\vq\cdot\vrho_3}\cdot
     e^{-i\vq'\cdot(\vrho_3-\vrho_4)}.
  \label{eq:v2_phases}
\end{align}
Multiplying and grouping:
\begin{equation}
  \text{(product of four)}
  = \underbrace{e^{+i\vq\cdot(\vrho_3-\vrho_2)}}_{F(\vq)}
    \;\cdot\;
    \underbrace{e^{-i\vq'\cdot\sS}}_{\text{$\vq'$-dep.}},
  \qquad
  \sS \equiv \vrho_1 - \vrho_2 + \vrho_3 - \vrho_4.
  \label{eq:phases_combined}
\end{equation}
Substituting into \cref{eq:sigma_phph} and pulling $F(\vq)$
outside the $\vq'$-sum:
\begin{align}
  \Sig(\omega;\vq)
  &= \frac{i\hbar}{2}
    \sum_{\substack{l_{1..4},\vrho_{1..4}\\j_{1..4}}}
    \Phi_1\,\Phi_2\;
    e^{+i\vq\cdot(\vrho_3-\vrho_2)}
    \int\!\dw
    \underbrace{\left[
      \frac{1}{N}\sum_{\vq'} e^{-i\vq'\cdot\sS}\;
      G_1(\omega';\vq')\;G_2(\omega{-}\omega';\vq{-}\vq')
    \right]}_{\displaystyle\text{convolution}},
  \label{eq:sigma_intermediate}
\end{align}
where $\Phi_1 \equiv \Phi^{ij_1j_2}_{l,l_1,l_2}$, etc.
The phase $e^{-i\vq'\cdot\sS}$ shifts the inverse DFT of $G_1$:
$g_1(\omega';\vR) \to g_1(\omega';\vR - \sS)$.
Applying the convolution theorem \cref{eq:conv_theorem}:
\begin{equation}
  \frac{1}{N}\sum_{\vq'} e^{-i\vq'\cdot\sS}\,
  G_1(\omega';\vq')\,G_2(\omega{-}\omega';\vq{-}\vq')
  = \sum_{\vR} g_1(\omega';\vR{-}\sS)\;
    g_2(\omega{-}\omega';\vR)\;e^{-i\vq\cdot\vR}.
  \label{eq:conv_result}
\end{equation}
Combining exponentials and substituting
$\vR' = \vR - (\vrho_3 - \vrho_2)$, the first Green's function
argument simplifies via
$\vrho_3 - \vrho_2 - \sS = \vrho_4 - \vrho_1$.
Incorporating \cref{eq:omega_conv}:
\begin{equation}
  \Sig(\omega;\vq)
  = \frac{i\hbar}{2}
    \sum_{\vR}
    \int\!d\tau\;
    \mathcal{B}(\tau;\vR)\,e^{+i\omega\tau}\,e^{-i\vq\cdot\vR},
  \label{eq:sigma_realspace}
\end{equation}
where
\begin{equation}
  \mathcal{B}^{ij}_{ll'}(\tau;\vR)
  \;\equiv\;
  \sum_{\substack{l_{1..4},\vrho_{1..4}\\j_{1..4}}}
  \Phi^{ij_1j_2}_{l,l_1,l_2}\;
  \Phi^{j_3j_4j}_{l',l_3,l_4}\;
  g^{j_1j_4}_{l_1 l_4}\!\bigl(\tau;\,\vR+\vrho_4{-}\vrho_1\bigr)\;
  g^{j_2j_3}_{l_2 l_3}\!\bigl(\tau;\,\vR+\vrho_3{-}\vrho_2\bigr).
  \label{eq:bubble_realspace}
\end{equation}
The $\vq'$-sum has been replaced by a loop over $N_D$ displacement
quadruples $(\vrho_1,\vrho_2,\vrho_3,\vrho_4)$.  At each
quadruple and each $\vR\in\mathcal{R}_\mathrm{eff}$, one inner
contraction \cref{eq:balanced_tree} is evaluated.  The total cost
is
\begin{equation}
  \cW_R
  = N_\tau\,N_D\,N_T^\mathrm{eff}\;\cC_\mathrm{ic},
  \qquad
  N_D \le N_\perp^4,\quad
  N_T^\mathrm{eff} = (4N_\mathrm{cut}+1)^2.
  \label{eq:cost_real}
\end{equation}

Although the contraction loop in \cref{eq:cost_real} is independent
of $N_q$, the displacement count $N_D \le N_\perp^4 =
(2N_\mathrm{cut}+1)^8$ grows as the eighth power of the FC3
cutoff.  The inner contraction $\cC_\mathrm{ic}$ is identical in
both formulations, so the real-space method is faster only when
$N_D\,N_T^\mathrm{eff} < N_q^2$.
For $N_q = 20\times 20 = 400$:
\begin{center}
\begin{tabular}{crr}
  \toprule
  $N_\mathrm{cut}$
    & $N_\perp^4 N_T^\mathrm{eff}$
    & ratio to $N_q^2$\\
  \midrule
  1 & $1.6\times 10^5$ & $1.0$  \\
  2 & $3.2\times 10^7$ & $198$ \\
  3 & $9.7\times 10^8$ & $6\,088$ \\
  \bottomrule
\end{tabular}
\end{center}
The crossover $N_q^2 = N_\perp^4 N_T^\mathrm{eff}$ occurs at
$N_q \approx 5{,}600$ for $N_\mathrm{cut}=2$, significantly larger than the
$N_q\approx 400$ required for phonons.

% ============================================================
\subsection{Separable FC3 approximation}
\label{sub:separable}
% ============================================================

Consider $\Phi^{i\mu_1\mu_2}_{l,l_1,l_2}$ as a function of the
transverse displacement pair $(\vrho_1,\vrho_2)$ at fixed slab and
combined Cartesian--basis indices.  Arranging these values into a
matrix with row index $(\vrho_1, \mu_1, l_1)$ and column index
$(\vrho_2, \mu_2, l_2)$, assume we can perform a separable
approximation
\begin{equation}
  \Phi^{i\mu_1\mu_2}_{l,l_1,l_2}
  \;\approx\;
  \sum_{r=1}^{R}
  f^{i\mu_1}_{l,l_1,r}(\vrho_1)\;h^{\mu_2}_{l_2,r}(\vrho_2),
  \label{eq:fc3_lowrank}
\end{equation}
with rank $R \ll N_\perp$, and similiarly
$\Phi^{\mu_3\mu_4 j}_{l',l_3,l_4} \approx \sum_s
  f'^{\mu_3 j}_{l',l_3,s}(\vrho_3)\;h'^{\mu_4}_{l_4,s}(\vrho_4)$
holds for the second FC3 factor.

Substituting \cref{eq:fc3_lowrank} for both FC3 factors in
\cref{eq:inner_contraction}:
\begin{align}
  \mathcal{B}^{ij}_{ll'}
  &\;\mathrel{+}=\;
  \sum_{r,s}
  \underbrace{
    \sum_{\mu_1,l_1,\mu_4,l_4}
    f^{i\mu_1}_{l,l_1,r}\;g^{\mu_1\mu_4}_{l_1,l_4}\;
    h'^{\mu_4}_{l_4,s}
  }_{P^i_{l,rs}}
  \;\cdot\;
  \underbrace{
    \sum_{\mu_2,l_2,\mu_3,l_3}
    h^{\mu_2}_{l_2,r}\;g^{\mu_2\mu_3}_{l_2,l_3}\;
    f'^{\mu_3 j}_{l',l_3,s}
  }_{Q^j_{l',rs}}.
  \label{eq:sep_split}
\end{align}
The four-tensor ring contraction has factorised into a product of
two independent bilinear forms $P$ and $Q$.  Each has the
structure $P_{rs} = f_r^{\,T}\,g_1\,h'_s$,  where
$f_r$ and $h'_s$ are vectors of dimension $M\hat{N}_c$ and $MN_B$
respectively and $g_1$ is the banded Green's function matrix.

The Green's function $g$ is shared across all $(r,s)$ pairs.
Rather than evaluating one bilinear form per $(r,s)$ pair at cost
$\cO(M^2\hat{N}_c N_B)$ each (totalling $R^2$ evaluations), we
can contract in two stages:
\begin{enumerate}
  \item\textbf{Matrix--vector products} (once per $s$):
    compute
    $\vv_s^{\mu_1}{}_{l_1}
      \equiv \sum_{\mu_4,l_4}
      g^{\mu_1\mu_4}_{l_1,l_4}\,h'^{\mu_4}_{l_4,s}$
    for $s = 1,\ldots,R$.
    Each product costs
    $\cO(M^2\hat{N}_c N_B)$, total $R\cdot M^2\hat{N}_c N_B$.
  \item\textbf{Inner products} (once per $(r,s)$):
    compute $P^i_{l,rs} = \sum_{\mu_1,l_1}
      f^{i\mu_1}_{l,l_1,r}\;\vv_s^{\mu_1}{}_{l_1}$,
    total $R^2\cdot M\hat{N}_c$.
\end{enumerate}
The combined cost per working point is
\begin{equation}
  2\bigl(R\cdot M^2\hat{N}_c N_B
        + R^2\cdot M\hat{N}_c\bigr).
  \label{eq:cost_opt}
\end{equation}
With $R \ll MN_B$, the first term dominates and the
effective cost is $\cO(R\,M^2\hat{N}_c N_B)$.

Fourier-transforming each factor of \cref{eq:fc3_lowrank},
\begin{equation}
  \hat{f}^{i\mu_1}_{l,l_1,r}(\vq)
  = \sum_{\vrho_1} f^{i\mu_1}_{l,l_1,r}(\vrho_1)\,
    e^{-i\vq\cdot\vrho_1},
  \qquad
  \hat{h}^{\mu_2}_{l_2,r}(\vq)
  = \sum_{\vrho_2} h^{\mu_2}_{l_2,r}(\vrho_2)\,
    e^{-i\vq\cdot\vrho_2},
  \label{eq:fh_fourier}
\end{equation}
the Fourier-transformed FC3 factorises as
$\PhiT(\vq_1,\vq_2) = \sum_r \hat{f}_r(\vq_1)\,\hat{h}_r(\vq_2)$.
Substituting into \cref{eq:sigma_phph} and using
\cref{eq:omega_conv}, the $\vq'$-sum for each $(r,s)$ pair
becomes
\begin{equation}
  \frac{1}{N}\sum_{\vq'}
  \underbrace{
    \hat{f}_r(\vq')\,\hat{h}'_s({-}\vq')\,g_1(\tau;\vq')
  }_{A_{rs}(\tau;\vq')}
  \;\underbrace{
    \hat{h}_r(\vq{-}\vq')\,\hat{f}'_s(\vq'{-}\vq)\,
    g_2(\tau;\vq{-}\vq')
  }_{B_{rs}(\tau;\vq{-}\vq')},
  \label{eq:lr_AB}
\end{equation}
a convolution evaluated via FFT.  Forming
$A_{rs}(\tau;\vq')$ is precisely the bilinear form
$\hat{f}_r^{\,T} G_1(\tau;\vq')\,\hat{h}'_s$, so the
optimal two-stage evaluation \cref{eq:cost_opt} applies:
first compute
$\vv_s(\tau;\vq') = G_1(\tau;\vq')\,\hat{h}'_s(\vq')$
for all $s$, then form
$A_{rs}(\tau;\vq') = \hat{f}_r(\vq')^T\,\vv_s(\tau;\vq')$
for all $(r,s)$.

The self-energy is
\begin{equation}
  \Sig(\omega;\vq)
  = \frac{i\hbar}{2}
    \int\!d\tau\;
    \sum_{r,s}
    \mathrm{FFT}_{\vR\to\vq}\!\bigl[
      a_{rs}(\tau;\vR)\;b_{rs}(\tau;\vR)
    \bigr]\;e^{+i\omega\tau},
  \label{eq:sigma_sep}
\end{equation}
where $a_{rs}, b_{rs}$ are the inverse DFTs of $A_{rs}, B_{rs}$.



% ======================================================================
\subsection{PCP decomposition of the FC3 tensor}
\label{sub:pcp}
% ======================================================================

The Permanent CP (PCP) decomposition, introduced in~\cite{luo_tensor_2025},
provides an alternative compressed
representation of the third-order force constants.  Unlike the SVD-based
separable approximation of \cref{sub:separable}, which decomposes the
FC3 viewed as a matrix (with rows and columns corresponding to the two
transverse partner atoms), the PCP decomposition works directly on the
three-index tensor and enforces full permutation symmetry of
the cubic interaction.

% -------------------------------------------------------
\subsubsection{PCP ansatz}
\label{sssub:pcp_ansatz}
% -------------------------------------------------------

The PCP ansatz approximates the mass-weighted FC3 tensor as
\begin{equation}\label{eq:pcp_ansatz}
  \tilde{\Psi}_{\alpha\beta\gamma}
    (0\kappa,\, l'\kappa',\, l''\kappa'')
  = \sum_{\xi=1}^{N_c}
    \frac{\lambda_\xi}{3!}
    \sum_{\sigma\in S_3}
    \sum_{\vc{R}}
    \prod_{i=1}^{3}
    A_{\sigma_i}^{\xi}\!\bigl(l_i - \vc{R},\, \kappa_i,\, \alpha_i\bigr),
\end{equation}
where:
\begin{itemize}
  \item $N_c$ is the PCP rank (number of components),
  \item $\lambda_\xi \in \mathbb{R}$ are the PCP singular values,
  \item $A_i^{\xi}(l, \kappa, \alpha) \in \mathbb{R}$ are the PCP mode
        functions, with $i \in \{1,2,3\}$ indexing three independent
        ``legs'' and $\xi$ indexing the rank component,
  \item the sum over $\sigma \in S_3$ (all 6 permutations of
        $\{1,2,3\}$) enforces full permutation symmetry,
  \item $\vc{R}$ runs over the $N_\mathrm{cell}$ unit cells in the
        supercell.
\end{itemize}

The PCP decomposition is structurally analogous to a canonical
polyadic (CP) decomposition of a third-order tensor, but with the
additional permanent (symmetrization) sum over $S_3$.  Each rank-$\xi$
component contributes a symmetrized outer product of three mode
vectors, weighted by $\lambda_\xi / 3!$.

The total number of parameters in the PCP ansatz is
\begin{equation}
  N_\mathrm{params}
  = 3 \, N_c \, N_\mathrm{cell} \, N_b \cdot 3 + N_c
  = N_c \,(9\, N_\mathrm{cell}\, N_b + 1),
\end{equation}
where the first term counts the mode entries (3 legs $\times$ $N_c$
ranks $\times$ $N_\mathrm{cell}$ cells $\times$ $N_b$ basis atoms
$\times$ 3 Cartesian directions) and the second term the singular
values.  For the Si $2\times 2\times 2$ supercell with
$N_b = 2$, $N_\mathrm{cell} = 8$, and $N_c = 8$:
$N_\mathrm{params} = 1{,}160$, compared to
$N_b \times N_\mathrm{super}^2 \times 27 = 13{,}824$ FC3 entries.
The compression ratio is $13{,}824 / 8 = 1{,}728\times$.

% -------------------------------------------------------
\subsubsection{Acoustic sum rule enforcement}
\label{sssub:pcp_asr}
% -------------------------------------------------------

A key advantage of the PCP representation is that the acoustic sum rule
(ASR) can be enforced directly on the mode functions.  Translational
invariance of the FC3 requires
\begin{equation}
  \sum_{\kappa} \Psi_{\alpha\beta\gamma}(\kappa, \kappa', \kappa'') = 0
  \qquad \forall\;(\kappa',\kappa'',\beta,\gamma).
\end{equation}
In the PCP parametrization, this reduces to the simpler condition
\begin{equation}\label{eq:pcp_asr}
  \sum_{l,\kappa} A_i^{\xi}(l, \kappa, \alpha) = 0
  \qquad \forall\;(\xi, i, \alpha),
\end{equation}
i.e., each mode function must have zero sum over all cell--atom pairs.
This is enforced during fitting via a penalty term
\begin{equation}
  \mathcal{L}_\mathrm{ASR}
  = \lambda_\mathrm{ASR}
    \sum_{\xi,i,\alpha}
    \left(
      \sum_{l,\kappa} A_i^{\xi}(l, \kappa, \alpha)
    \right)^{\!2},
\end{equation}
which is added to the reconstruction loss.  With a penalty weight
$\lambda_\mathrm{ASR} = 10$, the ASR violation
$\max|\sum_{l,\kappa} A|$ is driven below $10^{-4}$ during fitting.

This is in contrast to the separable SVD decomposition, where ASR
enforcement is more difficult.  The SVD factors $F_r$ and $H$ are
obtained from the stacked FC3 matrix restricted to same-slab atoms
($\delta l = 0$), and the partial sums over this restricted set of
atoms do not vanish even when the full ASR holds exactly.

The PCP mode functions and singular values are obtained by minimizing
the loss function
\begin{equation}\label{eq:pcp_loss}
  \mathcal{L}
  = \bigl\|\tilde{\Psi} - \Psi_\mathrm{target}\bigr\|_F^2
    + \mathcal{L}_\mathrm{ASR},
\end{equation}
where $\Psi_\mathrm{target}$ is the mass-weighted FC3 tensor (converted
to THz$^{5/2}$ units via \texttt{CONVERSION\_FC3\_THZ}) and
$\|\cdot\|_F$ is the Frobenius norm.

The optimization is implemented in PyTorch and uses a two-phase
strategy:
\begin{enumerate}
  \item \textbf{Adam warm-up} 500 iterations, learning
        rate 0.05 with cosine annealing to $10^{-3}$

  \item \textbf{L-BFGS refinement} remaining iterations,
        Hessian history size 50, strong Wolfe line search
\end{enumerate}

The forward model evaluation uses fully vectorized PyTorch operations.
The expensive computation is the triple outer product summed over cell shifts
and permutations:
\begin{equation}
  \tilde{\Psi}[\kappa_1, j, k, \alpha, \beta, \gamma]
  = \sum_{\xi}
    \frac{\lambda_\xi}{6}
    \sum_{\sigma \in S_3}
    \sum_l
    m_{\sigma_1}[\xi, \kappa_1, l, \alpha]\;
    m_{\sigma_2}[\xi, j, l, \beta]\;
    m_{\sigma_3}[\xi, k, l, \gamma],
\end{equation}
where the mode tables $m_i[\xi, j, l, \alpha]$ are precomputed via
advanced indexing from the cell-difference table.  This can be written as einsum
\begin{equation*}
  \texttt{fc3 += torch.einsum('xila, xjlb, xklc -> ijkabc', wm1, m2, m3)}
\end{equation*}
with $l$ (cell shift) and $x$ ($\xi$, rank) as contraction indices.

The gradient is computed by PyTorch automatic differentiation.  For the
Si test case ($N_c = 8$, 2000 iterations), fitting takes approximately
13\,s and achieves a relative Frobenius error of $4.4 \times 10^{-2}$.

\paragraph{Fitting results for Si $2\times2\times2$.}
\Cref{tab:pcp_fit} summarises the PCP fitting results for different
ranks.

\begin{table}[h]
\centering
\caption{PCP fitting results for Si $2\times2\times2$ FC3.}
\label{tab:pcp_fit}
\begin{tabular}{ccccc}
  \toprule
  $N_c$ & Parameters & Compression & Rel.\ error & ASR violation \\
  \midrule
  8  & 1{,}160  & 1{,}728$\times$ & $4.39\times10^{-2}$ & $5.3\times10^{-5}$ \\
  24 & 3{,}480  & 576$\times$     & $5.8\times10^{-3}$  & $2.7\times10^{-4}$ \\
  \bottomrule
\end{tabular}
\end{table}

To use the PCP decomposition in the $\vq$-space self-energy formulation,
the real-space mode functions are Fourier-transformed to each transverse
wavevector.  Define
\begin{equation}\label{eq:pcp_ft}
  f_i^{\xi}(\vq)_{\kappa\alpha}
  = \sum_l
    e^{-2\pi i\, \vq \cdot \vR_l}\;
    A_i^{\xi}(l, \kappa, \alpha),
\end{equation}
where $\vR_l$ is the lattice vector of cell $l$ in fractional
coordinates and $\vq$ is the transverse wavevector.  The result is a
complex vector of dimension $n_\mathrm{dof} = 3N_b$ for each
combination of leg index $i \in \{1,2,3\}$ and rank $\xi$.  No
additional mass weighting is applied since the mode functions were
fitted against the already mass-weighted target.

In matrix form, the Fourier-transformed modes at wavevector $\vq$ are
stored as
\begin{equation}
  f_\mathrm{all}(\vq) \in \mathbb{C}^{3 \times N_c \times n_\mathrm{dof}},
\end{equation}
with the three axes corresponding to leg index, rank, and combined
$(\kappa, \alpha)$ degree of freedom.

The ASR condition \cref{eq:pcp_asr} implies that at $\vq = \vc{0}$
(the $\Gamma$ point), the mode vectors satisfy
$\sum_{\kappa,\alpha} f_i^{\xi}(\vc{0})_{\kappa\alpha} \approx 0$,
suppressing the coupling to acoustic modes at long wavelengths.


% ======================================================================
\subsection{PCP self-energy kernel}
\label{sub:pcp_self_energy}
% ======================================================================

The PCP decomposition leads to a self-energy kernel with a qualitatively
different computational structure from both the dense and separable
formulations.  The FC3 is represented as a sum of symmetrized outer products,
and the bilinear contraction of two FC3 vertices with two Green's functions
can be reduced to scalar projections.

% -------------------------------------------------------
\subsubsection{Scalar projection}
\label{sssub:pcp_scalar_projection}
% -------------------------------------------------------

Substituting the PCP ansatz into the self-energy formula
\cref{eq:sigma_phph} and using the Fourier-transformed modes
\cref{eq:pcp_ft}, the self-energy becomes
\begin{align}\label{eq:sigma_pcp}
  \Sigma^{\gtrless}_{ab}(\omega; \vq)
  &= \frac{i\hbar}{2}
    \sum_{\xi,\xi'} \frac{\lambda_\xi \lambda_{\xi'}}{36}
    \sum_{\sigma,\sigma' \in S_3}
    \frac{1}{N} \sum_{\vq'}
    \int \frac{d\omega'}{2\pi}\;
    g_{\sigma_2 \sigma_2'}^{\xi \xi'}(\omega'; \vq')\;
    g_{\sigma_3 \sigma_3'}^{\xi \xi'}(\omega - \omega'; \vq - \vq')
    \notag\\
  &\qquad \times
    f_{\sigma_1}^{\xi}(\vq)_a \;
    \overline{f_{\sigma_1'}^{\xi'}(\vq)_b},
\end{align}
where the scalar projections of the Green's function are defined
as
\begin{equation}\label{eq:scalar_proj}
  g_{ij}^{\xi\xi'}(\omega; \vq)
  = \sum_{c,d}
    \overline{f_i^{\xi}(\vq)_c}\;
    G^{\gtrless}_{cd}(\omega; \vq)\;
    f_j^{\xi'}(\vq)_d
  = \vc{f}_i^{\xi}(\vq)^\dagger\;
    G^{\gtrless}(\omega; \vq)\;
    \vc{f}_j^{\xi'}(\vq).
\end{equation}

These are scalar-valued quantities: the $n_\mathrm{dof} \times
n_\mathrm{dof}$ matrix Green's function is projected onto the PCP mode
vectors.  The scalar
projections $g_{ij}^{\xi\xi'}$ carry indices
$(\xi, \xi', i, j)$ with $\xi, \xi' \in \{1, \ldots, N_c\}$ and
$i, j \in \{1, 2, 3\}$, giving $9 N_c^2$ scalar projections per
$(\omega, \vq)$ pair.

The frequency integral in \cref{eq:sigma_pcp} is a convolution of two
scalar functions, evaluated by FFT:
\begin{equation}
  c_{ij,kl}^{\xi\xi'}(\omega; \vq, \vq')
  = \mathrm{IFFT}\!\bigl[
    \hat{g}_{ij}^{\xi\xi'}(\vq') \cdot
    \hat{g}_{kl}^{\xi\xi'}(\vq - \vq')
  \bigr],
\end{equation}
where $\hat{g}$ denotes the FFT along the frequency axis.

% -------------------------------------------------------
\subsubsection{Accumulation and permutation batching}
\label{sssub:pcp_accumulation}
% -------------------------------------------------------

The self-energy is accumulated by iterating over external wavevector
pairs $(\vq, \vq')$ and permutation pairs $(\sigma, \sigma')$.  The
outer product structure of \cref{eq:sigma_pcp} allows the rank sum to
be vectorized.

For fixed $(\vq, \vq', \sigma, \sigma')$, the contribution to the
self-energy is
\begin{equation}\label{eq:pcp_accum}
  \Sigma^{\gtrless}_{ab}(\omega; \vq)
  \mathrel{+}=
  \frac{i\hbar \,\Delta\omega}{4\pi N}
  \sum_{\xi,\xi'}
  \frac{\lambda_\xi \lambda_{\xi'}}{36}\;
  c_{\sigma_2\sigma_2', \sigma_3\sigma_3'}^{\xi\xi'}(\omega)\;
  f_{\sigma_1}^{\xi}(\vq)_a\;
  \overline{f_{\sigma_1'}^{\xi'}(\vq)_b},
\end{equation}
which is computed as a single \texttt{einsum} operation:
\begin{equation*}
  \Sigma[\omega, a, b]
  \mathrel{+}=
  \mathrm{prefactor} \cdot
  \texttt{einsum}(\texttt{'xyw,xa,yb->wab'},\;
  W_{\xi\xi'} c^{\xi\xi'},\;
  f_{\sigma_1}^\xi,\;
  \overline{f_{\sigma_1'}^{\xi'}}),
\end{equation*}
where $W_{\xi\xi'} = \lambda_\xi \lambda_{\xi'} / 36$ is the weight
matrix.


% -------------------------------------------------------
\subsubsection{Computational cost}
\label{sssub:pcp_cost}
% -------------------------------------------------------

The PCP self-energy kernel has three phases per SCBA iteration:

\paragraph{Scalar projection (g-precomputation).}
For each wavevector $\vq$, compute $9N_c^2$ scalar projections from the
$n_\mathrm{dof} \times n_\mathrm{dof}$ Green's function matrices.
This is done as a batched matrix multiplication:
\begin{equation}
  G^{\gtrless}(\vq) \cdot \vc{f}^T
  \in \mathbb{C}^{n_\mathrm{fft} \times n_\mathrm{dof} \times 3N_c},
\end{equation}
followed by a contraction with $\overline{\vc{f}}$:
\begin{equation}
  g^{\xi\xi'}_{ij}[\vq, \omega]
  = \texttt{einsum}(\texttt{'ia,waj->ijw'},\;
    \overline{f}_\mathrm{flat},\; G\!\cdot\!f_\mathrm{flat}^T).
\end{equation}
Cost: $\mathcal{O}(N_q \cdot n_\mathrm{dof}^2 \cdot 3N_c \cdot
n_\mathrm{fft})$ for the matmul, plus
$\mathcal{O}(N_q \cdot n_\mathrm{dof} \cdot 9N_c^2 \cdot n_\mathrm{fft})$
for the contraction.

\paragraph{FFT of scalar projections.}
Batch FFT of the $g$ arrays along the frequency axis: cost
$\mathcal{O}(9N_c^2 \cdot N_q \cdot n_\mathrm{fft} \log n_\mathrm{fft})$.

\paragraph{Convolution and accumulation.}
For each $(\vq, \vq')$ pair: 36 element-wise products of
$(N_c, N_c, n_\mathrm{fft})$ arrays, grouped into a
$(N_c, N_c, 3, 3, n_\mathrm{fft})$ sum, one batch IFFT, and 9
rank-weighted einsums.  Cost per pair:
$\mathcal{O}(9N_c^2 n_\mathrm{fft} + N_c^2 \cdot 9 \cdot
n_\mathrm{fft} \log n_\mathrm{fft} + 9 N_c^2 n_\mathrm{freq}
n_\mathrm{dof})$.  Total over $N_q^2$ pairs:
\begin{equation}
  \mathcal{O}\!\left(
    N_q^2 \left[
      N_c^2 \cdot n_\mathrm{fft} \log n_\mathrm{fft}
      + N_c^2 \cdot n_\mathrm{freq} \cdot n_\mathrm{dof}
    \right]
  \right).
\end{equation}

The PCP kernel trades the $n_\mathrm{dof}^4$ scaling of the dense
formulation for a $N_c^2 n_\mathrm{dof}$ scaling in the accumulation
step, at the cost of introducing the fitting approximation error.


% ======================================================================
\subsection{Optimized separable self-energy}
\label{sub:separable_optimized}
% ======================================================================

The original separable self-energy kernel of \cref{sub:separable}
evaluates the $(\vq, \vq')$ convolution with an explicit loop over all
$R^2$ rank pairs $(r, s)$, performing 4 FFTs and 1 IFFT per pair.  For
the Si test case with $R = 24$ and $N_q = 16$ wavevectors, this results
in $R^2 \times 4 \times N_q^2 = 576 \times 4 \times 256 \approx
590{,}000$ individual FFT calls per SCBA iteration on a small
array of size $(n_\mathrm{fft}, n_\mathrm{dof})$.  The FFTs of the
intermediate arrays $V$ and $U$ can be precomputed once, and the
rank-pair sum can be factored analytically.  From \cref{sub:separable}:
\begin{equation}
  P_{rs}(\omega; \vq')
  = V_s(\omega; \vq') \cdot \hat{F}_r(\vq')^T,
  \qquad
  Q_{rs}(\omega; \vq - \vq')
  = U_r(\omega; \vq - \vq') \cdot \overline{\hat{F}_s(\vq - \vq')}^T,
\end{equation}
where $V_s = G \cdot \overline{\hat{h}_s}$ and $U_r = G^T \cdot
\hat{h}_r$ are the stage-1 products.

The self-energy contribution from a single $(\vq, \vq')$ pair is
\begin{equation}
  \Sigma_{ab}(\omega; \vq)
  \mathrel{+}=
  \mathrm{prefactor} \cdot
  \sum_{r,s}
  \mathrm{IFFT}\!\bigl[
    \hat{P}_{rs}(\cdot; \vq')_a \cdot
    \hat{Q}_{rs}(\cdot; \vq - \vq')_b
  \bigr](\omega).
\end{equation}

Since IFFT and summation commute, we can sum over $(r,s)$ in the
frequency domain before the inverse transform:
\begin{equation}
  \sum_{r,s} \hat{P}_{rs}[\omega, a] \cdot \hat{Q}_{rs}[\omega, b]
  = \sum_{c,d}
    \underbrace{
      \sum_r \hat{F}_r(\vq')_{ac} \cdot \hat{U}_r(\vq-\vq')_{\omega,d}
    }_{M_1[\omega,a,c,d]}
    \cdot
    \underbrace{
      \sum_s \hat{V}_s(\vq')_{\omega,c} \cdot
      \overline{\hat{F}_s(\vq-\vq')}_{bd}
    }_{M_2[\omega,c,b,d]},
\end{equation}
where $\hat{V}, \hat{U}$ are the FFTs of $V, U$ along the frequency
axis.  This gives three einsum:
\begin{align}
  M_1[\omega,a,c,d]
  &= \texttt{einsum}(\texttt{'rac,wdr->wacd'},\;
     \hat{F}(\vq'),\; \hat{U}(\vq{-}\vq')),
  \label{eq:M1} \\
  M_2[\omega,c,b,d]
  &= \texttt{einsum}(\texttt{'wcs,sbd->wcbd'},\;
     \hat{V}(\vq'),\; \overline{\hat{F}(\vq{-}\vq')}),
  \label{eq:M2} \\
  S[\omega,a,b]
  &= \texttt{einsum}(\texttt{'wacd,wcbd->wab'},\; M_1,\; M_2),
  \label{eq:S_contract}
\end{align}
followed by a single IFFT:
\begin{equation}
  \Sigma(\omega; \vq) \mathrel{+}=
  \mathrm{prefactor} \cdot \mathrm{IFFT}(S)[\mathrm{freq\_sl}].
\end{equation}

% ======================================================================
\subsection{Numerical validation}
\label{sub:validation_anharmonic}
% ======================================================================

The correctness of the self-energy kernels is validated with:
a synthetic toy model, a physical Si benchmark, and a comparison of all
three methods (dense, separable, PCP) on the same transport problem.

% -------------------------------------------------------
\subsubsection{Toy model validation}
\label{sssub:toy_model}
% -------------------------------------------------------

A synthetic test constructs random FC3 matrices $M_\mathrm{stacked}$
and gathering matrices $T(\vq)$ without any DFT data.  The separable
decomposition is obtained by SVD of $M_\mathrm{stacked}$, and the
Fourier-transformed factors $\hat{F}_r(\vq) = F_r \cdot T(\vq)^T$,
$\hat{H}(\vq) = T(\vq) \cdot H$ are constructed analytically.  Random
Hermitian Green's functions are generated at all $(\omega, \vq)$ points.

The following properties are verified:
\begin{enumerate}
  \item \textbf{Dense = Separable (full rank):}
    The self-energies $\Sigma^{<,>,R}$ computed by the dense kernel and
    the full-rank separable kernel agree to machine precision
    (relative difference $< 2 \times 10^{-15}$).

  \item \textbf{Rank truncation error:}
    Truncating the SVD to rank $R < R_\mathrm{max}$ introduces a
    self-energy error that decreases monotonically with increasing $R$
    and vanishes at full rank.

  \item \textbf{Self-energy relations:}
    The identity $\Sigma^R = \tfrac{1}{2}(\Sigma^> - \Sigma^<)$ holds
    to machine precision for all kernels.
\end{enumerate}


% -------------------------------------------------------
\subsubsection{Si benchmark results}
\label{sssub:si_benchmark}
% -------------------------------------------------------

The full comparison uses the Si $2\times2\times2$ supercell with FC3
from finite displacements.  All methods use a $4\times4$ transverse
$q$-mesh (16 points), a frequency grid of 101 points from 1.0 to
14.0\,THz, mixing parameter $\alpha = 0.3$, and 10 SCBA iterations.
The frequency range lower bound of 1.0\,THz avoids the acoustic mode
divergence at $\Gamma$, where
$G^<(\vq = \vc{0}, \omega \to 0) \propto n_\mathrm{BE}(\omega)/\omega^2$
diverges as $1/\omega^3$ for $\omega \to 0$.

\begin{table}[h]
\centering
\caption{Anharmonic thermal conductance for Si at 300\,K.  All methods
use identical harmonic inputs, frequency grid, and SCBA parameters.}
\label{tab:method_comparison}
\begin{tabular}{lccccc}
  \toprule
  Method & $G_\mathrm{ball}$ & $G_\mathrm{anh}$
    & $G_\mathrm{anh}/G_\mathrm{ball}$ & Conservation & Time \\
  & \multicolumn{2}{c}{[MW/(m$^2$\,K)]} & & & [s] \\
  \midrule
  Dense (exact)           & 1103 & 827 & 0.750 & $2.8\times10^{-2}$ & 53 \\
  Separable (full rank)   & 1103 & 827 & 0.750 & $2.8\times10^{-2}$ & 47 \\
  Separable ($R=6$)       & 1103 & 848 & 0.769 & $2.5\times10^{-2}$ & 22 \\
  PCP ($N_c = 8$)         & 1103 & 881 & 0.799 & $1.4\times10^{-2}$ & 91 \\
  PCP ($N_c = 24$)        & 1103 & 909 & 0.825 & $1.4\times10^{-2}$ & 791 \\
  \bottomrule
\end{tabular}
\end{table}

\begin{itemize}
  \item Dense and separable (full rank) agree to machine precision
    ($\Delta G / G < 10^{-15}$), confirming the mathematical
    equivalence demonstrated by the toy model.

  \item The truncated separable ($R = 6$) retains 81.5\% of the FC3
    Frobenius norm and overestimates $G_\mathrm{anh}$ by 2.5\%
    relative to the exact result, consistent with the reduced
    scattering strength from truncation.

  \item PCP overestimates $G_\mathrm{anh}$ by 6.5\% ($N_c = 8$) and
    9.9\% ($N_c = 24$).  The larger deviation at higher rank may
    indicate that PCP captures the dominant scattering channels
    differently from the slab-local SVD decomposition, though the
    fitting error of 4.4\% ($N_c = 8$) and 0.6\% ($N_c = 24$) in the
    FC3 reconstruction does not directly translate to a proportional
    transport error.

  \item PCP consistently achieves better heat flow conservation
    ($\sim 1.4\%$ vs.\ $\sim 2.5$--$2.8\%$ for dense/separable)

  \item The frequency lower bound of 1.0\,THz is essential for
    convergence of the dense and separable methods.  With
    $\omega_\mathrm{min} = 0.01$\,THz, the SCBA with full FC3 diverges due to
    the acoustic-mode singularity at $\Gamma$.  PCP converges even at
    $\omega_\mathrm{min} = 0.01$\,THz
\end{itemize}


% ======================================================================
\subsection{I/O, Checkpoints}
\label{sub:io}
% ======================================================================

To enable reproducible comparisons and avoid redundant recomputation,
the transport results and FC3 decompositions are saved to disk using
a standardised format implemented in \texttt{phonon\_inputs/io.py}.

Transport results are stored as a pair of files:
\begin{itemize}
  \item \texttt{transport.npz}: compressed NumPy archive containing
    array data (spectral heat currents, self-energies, transmission,
    frequency grid),
  \item \texttt{transport.json}: JSON sidecar with scalar quantities
    (thermal conductances, conservation errors, iteration counts) and
    optional metadata (method name, parameters, timing).
\end{itemize}

The functions \texttt{save\_transport\_results} and
\texttt{load\_transport\_results} handle serialisation and
deserialisation, including automatic type dispatch (arrays to NPZ,
scalars to JSON).  An analogous pair
\texttt{save\_decomposition}/\texttt{load\_decomposition} stores SVD or
PCP decompositions for reuse across SCBA runs.

The comparison script \texttt{compare\_methods.py} runs all five method
variants (dense, separable full rank, separable $R=6$, PCP $N_c=8$, PCP
$N_c=24$), saves each result to a subdirectory of
\texttt{anharmonic/results/}, and generates a multi-panel diagnostic
plot.  Subsequent runs with the \texttt{-{}-load} flag skip the
computation and regenerate plots from the saved data, enabling rapid
iteration on visualisation without rerunning the SCBA.

% ======================================================================
\subsection{Decomposition Results}
\label{sub:decomp_results}
% ======================================================================

We first present the SVD approximation quality
(\cref{sssub:pcp_quality}), cost analysis (\cref{sssub:cost_analysis}),
and thickness-dependent validation (\cref{sssub:thickness_validation}).
\Cref{sssub:scba_investigation} investigates SCBA convergence limits
using a toy model and FC3/FC2 strength analysis.
\Cref{sssub:svd_asr} evaluates ASR enforcement.


% -------------------------------------------------------
\subsubsection{Approximation quality}
\label{sssub:pcp_quality}
% -------------------------------------------------------

All results in this section use the Si $2\times 2\times 2$ test case
with a $4\times 4$ $q$-mesh, 101 frequency points from
1.0--14.0\,THz, up to 10 SCBA iterations (tol.\ 0.5\%), and mixing
0.3.  The dense reference gives
$G_\mathrm{anh} = 827\,\mathrm{MW/(m^2\,K)}$.

\paragraph{FC3 reconstruction error.}
\Cref{fig:fc3_error} shows the singular value spectrum of the
stacked FC3 matrix (left) and the reconstruction error vs.\ SVD
rank (right).  The error decreases monotonically by construction,
reaching $< 3\%$ at $R = 12$ and machine precision at $R = 24$
(the full rank).

\begin{figure}[htbp]
  \centering
  \includegraphics[width=\textwidth]{fig/decomp/fc3_error_vs_rank.pdf}
  \caption{(a) Singular value spectrum of the stacked FC3 matrix.
    (b) FC3 reconstruction error vs.\ SVD rank.}
  \label{fig:fc3_error}
\end{figure}

\paragraph{Transport accuracy vs.\ rank.}
\Cref{fig:transport_rank} shows $G_\mathrm{anh}$ and the relative
transport error as a function of SVD rank.  The separable
approximation converges monotonically to the dense result, reaching
$< 1\%$ error at $R = 12$ and matching dense exactly at $R = 24$.

\begin{figure}[htbp]
  \centering
  \includegraphics[width=\textwidth]{fig/decomp/transport_vs_rank.pdf}
  \caption{(a) Thermal conductance $G_\mathrm{anh}$ vs.\ SVD rank.
    The separable result (green) converges to the dense reference
    (red line).
    (b) Relative transport error.  The green band marks the $\pm 5\%$
    region; SVD enters this band at $R \geq 6$.}
  \label{fig:transport_rank}
\end{figure}

\begin{table}[h]
\centering
\caption{Transport results for Si 300\,K.  Dense reference:
$G_\mathrm{anh} = 827\,\mathrm{MW/(m^2\,K)}$.}
\label{tab:accuracy}
\begin{tabular}{lccccc}
  \toprule
  Method & FC3 err. & $G_\mathrm{anh}$ [MW] & Transp.\ err.
         & Conserv. & Time [s] \\
  \midrule
  Dense           & ---     & 827  & ---       & 2.8\% & 53 \\
  SVD $R=2$       & 43.3\%  & 1054 & $+27.5\%$ & 0.4\% & 4 \\
  SVD $R=4$       & 22.2\%  & 957  & $+15.7\%$ & 1.2\% & 20 \\
  SVD $R=6$       & 12.3\%  & 848  & $+2.5\%$  & 2.5\% & 22 \\
  SVD $R=12$      & 2.8\%   & 831  & $+0.4\%$  & 2.8\% & 31 \\
  SVD $R=24$      & 0.1\%   & 827  & $-0.0\%$  & 2.8\% & 47 \\
  \bottomrule
\end{tabular}
\end{table}

\paragraph{Spectral heat current.}
\Cref{fig:spectral_current} shows the frequency-resolved heat
current (left) and its deviation from the dense reference (right).
SVD $R = 24$ is indistinguishable from dense.

\begin{figure}[htbp]
  \centering
  \includegraphics[width=\textwidth]{fig/decomp/spectral_current.pdf}
  \caption{(a) Spectral heat current for selected methods.
    (b) Difference from dense: SVD $R = 24$ is indistinguishable from
    dense; SVD $R = 6$ shows small deviations.}
  \label{fig:spectral_current}
\end{figure}

\paragraph{Self-energy comparison.}
\Cref{fig:sigma_compare} shows the phonon broadening
$-\mathrm{Tr}[\mathrm{Im}\,\Sigma^R]/M$ (left) and the relative
Frobenius error of the self-energy vs.\ frequency (right).
SVD $R = 24$ reproduces the dense self-energy to $< 1\%$ at all
frequencies.

\begin{figure}[htbp]
  \centering
  \includegraphics[width=\textwidth]{fig/decomp/self_energy_comparison.pdf}
  \caption{(a) Phonon broadening (q-averaged).
    (b) Relative Frobenius error of the self-energy vs.\ frequency.
    SVD converges systematically with increasing rank.}
  \label{fig:sigma_compare}
\end{figure}

\paragraph{FC3 error vs.\ transport error.}
\Cref{fig:error_decomp} plots the FC3 reconstruction error against
the transport error.  SVD points follow the
$\epsilon_G \approx C \delta_\Phi$ line: transport error is directly
proportional to the FC3 truncation error.

\begin{figure}[htbp]
  \centering
  \includegraphics[width=0.7\textwidth]{fig/decomp/error_decomposition.pdf}
  \caption{FC3 reconstruction error vs.\ transport error.
    SVD (green squares) follows $\epsilon_G \propto \delta_\Phi$:
    better FC3 $\Rightarrow$ better transport.}
  \label{fig:error_decomp}
\end{figure}

\paragraph{Conservation and cost trade-off.}
\Cref{fig:conserv_pareto} shows conservation error (left) and the
accuracy--cost Pareto front (right).  SVD $R = 6$ (22\,s, 2.5\%
error) and SVD $R = 12$ (31\,s, 0.4\% error) provide excellent
accuracy--cost trade-offs.

\begin{figure}[htbp]
  \centering
  \includegraphics[width=\textwidth]{fig/decomp/conservation_and_timing.pdf}
  \caption{(a) Heat flow conservation error.
    (b) Pareto front: $G_\mathrm{anh}$ accuracy vs.\ wall time.
    The green band marks $\pm 5\%$ of the dense reference.}
  \label{fig:conserv_pareto}
\end{figure}

\paragraph{SCBA convergence.}
\Cref{fig:scba_conv} shows the SCBA convergence behaviour.  All
SVD ranks converge at similar rates, requiring $\sim$10 iterations.
Low-rank cases ($R = 2$) converge faster.

\begin{figure}[htbp]
  \centering
  \includegraphics[width=\textwidth]{fig/decomp/scba_convergence.pdf}
  \caption{(a) SCBA convergence: relative change in total heat
    current vs.\ iteration.
    (b) Total iterations to convergence.}
  \label{fig:scba_conv}
\end{figure}


% -------------------------------------------------------
\subsubsection{Cost analysis}
\label{sssub:cost_analysis}
% -------------------------------------------------------

The separable kernel replaces the dense per-pair cost with a sum
over $R^2$ rank pairs.  For each $(\vq,\vq')$ pair:
\begin{equation}\label{eq:sep_cost}
  \cW_\mathrm{sep}^{\mathrm{pair}}
  = \underbrace{R^2\,N_\tau \log N_\tau}_{\text{batch IFFT}}
  + \underbrace{R^2\,N_\omega\,M^2}_{\text{accumulation}},
\end{equation}
where $M = n_\mathrm{dof} = 6$ for Si.  The dense kernel has cost
\begin{equation}
  \cW_\mathrm{dense}^{\mathrm{pair}}
  = M^4\,N_\tau \log N_\tau + N_\omega\,(M^5 + M^4).
\end{equation}
The separable kernel is faster when $R^2 < M^4$, i.e., $R < M^2$.
For Si ($M = 6$): at $R = 6$, $R^2 = 36 \ll M^4 = 1296$, giving
$\sim 36\times$ fewer IFFT operations.


% -------------------------------------------------------
\subsubsection{Thickness-dependent thermal conductivity}
\label{sssub:thickness_validation}
% -------------------------------------------------------

To validate the separable kernel,
we reproduce the thickness-dependent cross-plane thermal conductivity
of a silicon thin film at 300\,K, following \cite{guo_quantum_2020}.  The setup models a thin film of thickness~$d$
between two semi-infinite Si contacts, with a temperature difference
$\Delta T = 10$\,K.  The thermal conductivity is computed as
$\kappa_\mathrm{eff} = G \cdot d$, where $G = J / (A_c \Delta T)$ is
the thermal conductance and $d = n_\mathrm{slabs} \times a/2$ is the
film thickness (one FCC primitive cell per slab, $a/2 = 2.73$\,\AA{}
along~$[100]$).

\Cref{fig:thickness_kappa} shows the thickness-dependent thermal
conductivity computed with the separable kernel ($R = 12$, $4 \times 4$
$q$-mesh, 101 frequency points, 0.5--15.5\,THz).  Both the ballistic
and anharmonic results are compared with the data of
Guo et~al.\ (first nearest-neighbor FC3, DFT input).

\begin{figure}[htbp]
  \centering
  \includegraphics[width=\textwidth]{fig/decomp/thickness_conductivity.pdf}
  \caption{(a) Cross-plane thermal conductivity of a Si thin film at
    300\,K vs.\ thickness, computed with the separable kernel
    ($R = 12$) and compared with Guo et~al.\ \cite{guo_quantum_2020}.
    Open symbols: reference data; filled symbols: present work.
    (b) Thermal conductance vs.\ thickness, showing the ballistic
    (circles) and anharmonic (triangles) results.}
  \label{fig:thickness_kappa}
\end{figure}

The anharmonic reduction of thermal conductivity should increase with film
thickness: thicker films allow more phonon-phonon
scattering.  The SCBA converges reliably up to $n_\mathrm{slabs} = 6$
($d \approx 1.6$\,nm).  At $n_\mathrm{slabs} = 10$, the solver yields 47.8\% conservation error.
For the converged data points ($n_\mathrm{slabs} \leq 6$), the
anharmonic reduction is 30--50\%, qualitatively consistent with
Guo et~al.\ but quantitatively larger.

\Cref{fig:thickness_spectral} shows the spectral heat current at
different thicknesses, there are still problems..

\begin{figure}[htbp]
  \centering
  \includegraphics[width=\textwidth]{/decomp/thickness_spectral.pdf}
  \caption{(a) Spectral heat current vs.\ frequency for different
    film thicknesses (solid: anharmonic, dashed: ballistic).
    (b) Anharmonic conductance reduction vs.\ thickness.}
  \label{fig:thickness_spectral}
\end{figure}


% -------------------------------------------------------
\subsubsection{SCBA convergence investigation}
\label{sssub:scba_investigation}
% -------------------------------------------------------

The thickness sweep of \cref{sssub:thickness_validation} reveals that
the SCBA diverges for thicknesses above $\sim$6 primitive slabs
($\gtrsim 1.6$\,nm), even with reduced mixing parameters.  This
section investigates the root cause by (a) validating the solver
equations with a toy model, (b) quantifying the anharmonic coupling
strength ratio FC3/FC2, and (c) identifying the physical regime in
which the perturbative SCBA is valid.

To distinguish solver bugs from physical convergence limits, we
construct a synthetic simple-cubic crystal with one atom per cell
($a = 3$\,\AA, $m = 28$\,amu), nearest-neighbor harmonic springs
($k = 5$\,eV/\AA$^2$), and nearest-neighbor cubic anharmonicity
($\Phi_3$, isotropic, fully $S_3$-symmetrized).  This model is fed
through the \emph{existing} separable SCBA code.  \Cref{tab:toy_weak} shows that for weak anharmonicity
($\Phi_3 = 0.5$\,eV/\AA$^3$), the SCBA converges at all tested device
thicknesses ($n_\mathrm{slabs} = 1$--$10$) with heat-flow conservation
$<$3.8\% and physical conductance reduction with increasing thickness.

\begin{table}[h]
  \centering
  \caption{SCBA convergence for the toy model with weak anharmonicity
    ($\Phi_3 = 0.5$\,eV/\AA$^3$, $\alpha = 0.001$).
    All runs use $2\times 2$ $q$-mesh and mixing~0.3.}
  \label{tab:toy_weak}
  \begin{tabular}{rrrrrl}
    \toprule
    $n_\mathrm{slabs}$ & $G_\mathrm{ball}$ & $G_\mathrm{anh}$
      & Conserv.\ & Iters & Converged \\
    & (MW/m$^2$K) & (MW/m$^2$K) & & & \\
    \midrule
    1  & 1302 & 1297 & 0.56\% & 40 & yes \\
    2  & 1294 & 1286 & 1.05\% & 40 & yes \\
    4  & 1279 & 1264 & 1.87\% & 40 & yes \\
    6  & 1264 & 1244 & 2.57\% & 40 & yes \\
    8  & 1250 & 1224 & 3.19\% & 40 & yes \\
    10 & 1236 & 1206 & 3.73\% & 40 & yes \\
    \bottomrule
  \end{tabular}
\end{table}

Increasing the coupling to $\Phi_3 = 5.0$\,eV/\AA$^3$ causes the
SCBA to diverge even at $n_\mathrm{slabs} = 1$, regardless of mixing
parameter (0.3, 0.1, or 0.05) or $q$-mesh ($2\times 2$ or
$4\times 4$).  \Cref{tab:toy_strong} shows results:
conductance exceeds the ballistic limit and/or conservation errors
exceed 5\%. One configuration appear marginally stable
($n_\mathrm{slabs} = 2$, $q$-mesh $2\times 2$, mixing 0.1).  At $n_\mathrm{slabs} = 10$--$20$, the
self-energy blows up to $\sim$$10^4\times$ the ballistic conductance.

\begin{table}[h]
  \centering
  \caption{SCBA results for strong anharmonicity
    ($\Phi_3 = 5.0$\,eV/\AA$^3$, $\alpha = 0.011$)}
  \label{tab:toy_strong}
  \begin{tabular}{rrrlrrrl}
    \toprule
    $n_\mathrm{sl}$ & $q$-mesh & mix & $G_\mathrm{ball}$
      & $G_\mathrm{anh}$ & Conserv. & Iters & OK \\
    \midrule
    1 & $4\times 4$ & 0.30 & 1302 & 1105 & 18.8\% & 40 & no \\
    1 & $2\times 2$ & 0.05 & 1302 & 1164 & 10.2\% & 40 & no \\
    \textbf{2} & \textbf{2$\times$2} & \textbf{0.10}
      & \textbf{1294} & \textbf{1117} & \textbf{1.5\%}
      & \textbf{40} & \textbf{yes} \\
    4 & $4\times 4$ & 0.10 & 1279 &  884 & 9.7\%  & 40 & no \\
    10 & $2\times 2$ & 0.30 & 1236 & 15585 & 100\% & 40 & no \\
    20 & $4\times 4$ & 0.10 & 1169 & 26313 & 100\% & 40 & no \\
    \bottomrule
  \end{tabular}
\end{table}

A larger $q$-mesh makes convergence harder.  At $n_\mathrm{slabs} = 1$ with $\Phi_3 = 5.0$, the
conservation error is 7\% for $2\times 2$ vs.\ 19\% for
$4\times 4$ at the same mixing.  More $q$-points introduce more
coupled self-energy channels, each contributing to the feedback loop.

\paragraph{FC3/FC2 analysis: Si vs.\ toy model.}
\Cref{tab:alpha} compares the anharmonic parameter for the real Si
$2\times 2\times 2$ FC3 and the toy model at different $\Phi_3$
strengths.

We define a dimensionless anharmonic parameter
\begin{equation}\label{eq:alpha}
  \alpha = \frac{\Phi_{3,\max}\,u_\mathrm{rms}}{\Phi_{2,\mathrm{diag}}},
\end{equation}
where $\Phi_{3,\max}$ is the maximum Frobenius norm of any FC3 triplet
block, $\Phi_{2,\mathrm{diag}}$ is the mean on-site FC2 norm, and
$u_\mathrm{rms} = \sqrt{k_B T / \Phi_{2,\mathrm{diag}}}$ is the
thermal displacement amplitude.  The parameter $\alpha$ represents the
ratio of anharmonic to harmonic force at the thermal displacement
scale.

\begin{table}[h]
  \centering
  \caption{Anharmonic parameter $\alpha$ for Si and the toy model at
    300\,K.}
  \label{tab:alpha}
  \begin{tabular}{lrrrrl}
    \toprule
    System & $\Phi_{2,\mathrm{diag}}$ & $\Phi_{3,\max}$
      & $u_\mathrm{rms}$ (\AA) & $\alpha$ & SCBA \\
    & (eV/\AA$^2$) & (eV/\AA$^3$) & & & \\
    \midrule
    Si (phono3py) & 23.1 & 74.8 & 0.033 & 0.108 & marginal \\
    Toy, $\Phi_3 = 0.5$ & 8.66 & 0.17 & 0.055 & 0.001 & converges \\
    Toy, $\Phi_3 = 2.0$ & 8.66 & 0.67 & 0.055 & 0.004 & fails \\
    Toy, $\Phi_3 = 5.0$ & 8.66 & 1.67 & 0.055 & 0.011 & fails \\
    \bottomrule
  \end{tabular}
\end{table}

\paragraph{Comparison.}
\cite{guo_quantum_2020} successfully computes
$\kappa$ up to $d = 13$\,nm (24 conventional unit cells = 48 primitive
slabs). Two differences:
\begin{enumerate}
  \item \textbf{Larger FC3 supercell}: $3\times 3\times 3$ (216~atoms)
    vs.\ our $2\times 2\times 2$ (16~atoms).
  \item \textbf{First-NN-only FC3}: Their Table~IV analysis uses only
    first-nearest-neighbor FC3, which filters out the
    longer-range contributions that increase $\Phi_{3,\max}$.
\end{enumerate}
I should try
\begin{enumerate}
  \item \textbf{Better FC3 input}: FC3 from a $3\times 3\times 3$
    supercell with controlled neighbor-shell truncation to ensure
    $\alpha \lesssim 0.01$.
  \item \textbf{Different mixing}: Replace simple linear mixing
    $\Sigma_\mathrm{new} = (1-\alpha_\mathrm{mix})\Sigma_\mathrm{old}
    + \alpha_\mathrm{mix}\Sigma_\mathrm{computed}$ with Anderson or
    DIIS acceleration, dynamically reduce
    $\alpha_\mathrm{mix}$ when $|\Sigma|$ grows.
\end{enumerate}

% ============================================================
\subsection{Computational scaling}
\label{sub:scaling}
% ============================================================

\subsubsection{Cost summary}

The dominant cost of each formulation is:
\begin{align}
  \cW_q &= N_\tau\,N_q^2\;\cC_\mathrm{ic},
  &\cC_\mathrm{ic} &= \cO(M^3\hat{N}_c^2 N_B),
  \tag{$\vq$-space}
  \label{eq:scaling_q}\\[3pt]
  \cW_R &= N_\tau\,N_D\,N_T^\mathrm{eff}\;\cC_\mathrm{ic},
  &N_D &\le N_\perp^4,
  \tag{real space}
  \label{eq:scaling_R}\\[3pt]
  \cW_S &= N_\tau\,N_q\bigl[
      R\,M^2\hat{N}_c N_B + R^2 M\hat{N}_c
    \bigr],
  && \tag{separable}
  \label{eq:scaling_S}
\end{align}
plus $\cO(R^2 N_\tau N_q\log N_q\,M)$ for the FFT
convolutions.  The $\vq$-space and real-space formulations share
the same inner contraction $\cC_\mathrm{ic}$ and differ only in
their outer loop count ($N_q^2$ vs.\ $N_D N_T^\mathrm{eff}$).
The separable formulation differs in both respects: the outer loop
is $R^2 N_q$ convolutions (each via FFT at $N_q\log N_q$) and the
inner operation is a bilinear form rather than a ring contraction.

In the regime $R \ll MN_B$ where the matrix--vector products
dominate, the speedup relative to $\vq$-space is
\begin{equation}
  \frac{\cW_q}{\cW_S}
  \;\approx\;
  \frac{N_q\,M\,\hat{N}_c}{R}.
  \label{eq:speedup}
\end{equation}
This ratio grows linearly in both $N_q$ and the cross-section size
$M = CN_s$, and inversely in $R$.

\subsubsection{Memory requirements}

\begin{center}
\begin{tabular}{lll}
  \toprule
  Formulation & Dominant array & Size \\
  \midrule
  $\vq$-space
    & $g(\tau;\vq)$ for all $(\tau,\vq)$
    & $N_\tau\,N_q\,M^2 N_B$ \\
    & FC3 (real space, on-the-fly)
    & $N_D\,M^2\hat{N}_c^2$ \\[4pt]
  Real space
    & $g(\tau;\vR)$ for $\vR\in\mathcal{R}_\mathrm{eff}$
    & $N_\tau\,N_T^\mathrm{eff}\,M^2 N_B$ \\
    & FC3 (real space)
    & $N_D\,M^2\hat{N}_c^2$ \\[4pt]
  Separable
    & $\mathcal{B}(\tau;\vq)$ (accumulated bubble)
    & $N_\tau\,N_q\,M^2 N_B$ \\
    & $G(\tau;\vq)$ for current $\tau$
    & $N_q\,M^2 N_B$ \\
    & Rank factors $\hat{f}_r(\vq), \hat{h}_r(\vq)$
    & $R\,N_q\,M\hat{N}_c$ \\
    & $\vv_s(\vq)$, $A_{rs}(\vq)$, $B_{rs}(\vq)$
    & $R\,N_q\,M\hat{N}_c + 2R^2\,N_q$ \\
  \bottomrule
\end{tabular}
\end{center}

% ============================================================
\subsection{Distributed computation of the separable
  phonon--phonon self-energy}
\label{sub:distributed}
% ============================================================

We describe the distributed algorithm for computing the
separable phonon--phonon self-energy
\cref{eq:sigma_sep} within a large-scale NEGF solver.
The algorithm builds on the existing parallelisation
infrastructure of QuaTrEx~\cite{vetsch_ab-initio_2025}, which
distributes data over energy points and spatial partitions,
and extends it with a transverse-wavevector ($\vq$) dimension
required for the phonon--phonon interaction.

% -------------------------------------------------------
\subsubsection{Rank layout}
\label{sssub:rank_layout}
% -------------------------------------------------------

The total number of MPI ranks $P$ is factored as
\begin{equation}
  P = P_\omega \times P_q \times P_S,
  \label{eq:rank_layout}
\end{equation}
where $P_\omega$ ranks share the $N_\omega$ frequency points,
$P_q$ ranks share the $N_q$ transverse wavevectors, and $P_S$
ranks participate in the spatial domain decomposition of the
RGF solver along the transport direction.  Each rank is
identified by a triple $(p_\omega, p_q, p_S)$.  The
$P_\omega$-ranks at fixed $(p_q, p_S)$ form an
$\omega$-communicator, the $P_q$-ranks at fixed
$(p_\omega, p_S)$ form a $\vq$-communicator, and the
$P_S$-ranks at fixed $(p_\omega, p_q)$ form a
spatial communicator used for the distributed RGF.

In the RGF phase (solving for $G$), each rank stores the
diagonal and first off-diagonal blocks of
$G^{\gtrless}(\omega_k;\vq_n)$ for its assigned
$(\omega_k,\vq_n)$-points, with block size $N_{BS}\times N_{BS}$.
The number of frequency points per $\omega$-rank is
$N_\omega^{\mathrm{loc}} = N_\omega / P_\omega$, and the number
of $\vq$-points per $\vq$-rank is
$N_q^{\mathrm{loc}} = N_q / P_q$.  For the spatial
decomposition with $P_S$ partitions, each spatial rank holds
$\sim N_B/P_S$ transport cells of the block-tridiagonal system.

% -------------------------------------------------------
\subsubsection{Precomputation of rank factors}
\label{sssub:precompute}
% -------------------------------------------------------

The separable FC3 approximation \cref{eq:fc3_lowrank} is
computed once before the SCBA loop.  The real-space rank
factors $f^{i\mu_1}_{l,l_1,r}(\vrho_1)$ and
$h^{\mu_2}_{l_2,r}(\vrho_2)$ are obtained by truncated SVD
of the FC3 matrix (viewed as a matrix in $(\vrho_1,\mu_1,l_1)
\times (\vrho_2,\mu_2,l_2)$) to rank $R$.
Their Fourier transforms $\hat{f}_r(\vq)$, $\hat{h}_r(\vq)$ are
computed by applying 2D DFTs over the $N_\perp$ displacement
vectors for each $(r, \mu, l)$-combination.  Since the rank
factors have dimension $M\hat{N}_c$ per slab group and there
are $R$ of them, the total storage is
$\cO(R\,N_q\,M\hat{N}_c)$ per FC3 factor.  This is small
compared to the Green's function storage and is replicated on
all $\vq$-ranks.

% -------------------------------------------------------
\subsubsection{Algorithm overview}
\label{sssub:algorithm}
% -------------------------------------------------------

Each SCBA iteration for the phonon--phonon self-energy
proceeds through five phases, detailed below.  We write $G_k^n \equiv
G^{\gtrless}(\omega_k;\vq_n)$ for the Green's function at
frequency $\omega_k$ and wavevector $\vq_n$.

\paragraph{Phase~1: $\omega\to\tau$ transform
  (all-to-all within $\omega$-communicators).}
%
The RGF solver produces $G_k^n$ distributed over
$(\omega,\vq)$-pairs: each rank $(p_\omega,p_q,p_S)$ holds
$N_\omega^{\mathrm{loc}} \times N_q^{\mathrm{loc}}$ points.
To apply the inverse FFT $\omega\to\tau$, all $N_\omega$
frequencies at a given $(\vq_n, \text{entry})$ must be
gathered on a single rank.

This is done by the Stacked$\,\to\,$NNZ transpose:
an all-to-all exchange within each $\omega$-communicator
redistributes the data so that each rank holds all $N_\omega$
frequencies for a subset of $(\text{entry},\vq)$-pairs.
After the transpose, each rank applies the $\omega\to\tau$
IFFT independently.  The result is
$g^{\mu_1\mu_4}_{l_1,l_4}(\tau_m;\vq_n)$ for all $N_\tau$
time points, at the rank's assigned $(\text{entry},\vq)$-pairs.

The communication volume per rank is
$N_\omega^{\mathrm{loc}}\times N_q^{\mathrm{loc}} \times
  N_{\mathrm{nnz}}^{\mathrm{loc}}$ complex numbers, where
$N_{\mathrm{nnz}}^{\mathrm{loc}}$ is the number of nonzero
Green's function entries per $(\omega,\vq)$-point assigned to
this rank (diagonal and first off-diagonal blocks of the
block-tridiagonal structure, divided over $P_S$ spatial
partitions).

\paragraph{Phase~2: bilinear contraction (local).}
%
This is computationally dominant.
At each $(\tau_m, \vq_n)$ assigned to the current rank, the
two-stage contraction is executed:

\begin{enumerate}
\item For $s = 1,\ldots,R$: compute the matrix--vector product
  \begin{equation}
    \vv_s^{\mu_1}{}_{l_1}(\tau_m;\vq_n)
    = \sum_{\mu_4,l_4}
      g^{\mu_1\mu_4}_{l_1,l_4}(\tau_m;\vq_n)\;
      \hat{h}'^{\mu_4}_{l_4,s}(\vq_n).
    \label{eq:matvec}
  \end{equation}
  The Green's function $g(\tau_m;\vq_n)$ is block-tridiagonal
  with $N_B/P_S$ blocks of size $N_{BS}\times N_{BS}$ on the
  current spatial partition.  The vector
  $\hat{h}'_s(\vq_n)$ is nonzero only within $\hat{N}_c$ slabs
  of the output slab $l'$, so the product involves at most
  $\hat{N}_c$ block-rows of $g$, each with 3 contributing
  blocks (block-tridiagonal structure).  The cost per
  $(s,\tau_m,\vq_n)$ is $\cO(\hat{N}_c\,N_{BS}^2)$.

  The result $\vv_s$ is a vector of dimension $M\hat{N}_c$ per
  output slab, and is reused across all $R$ values of $r$ in
  the next step.

\item For $r = 1,\ldots,R$ and $s = 1,\ldots,R$: compute the
  inner product
  \begin{equation}
    A^{i}_{l,rs}(\tau_m;\vq_n)
    = \sum_{\mu_1,l_1}
      \hat{f}^{i\mu_1}_{l,l_1,r}(\vq_n)\;
      \vv_s^{\mu_1}{}_{l_1}(\tau_m;\vq_n).
    \label{eq:innerprod}
  \end{equation}
  Each inner product has length $M\hat{N}_c$.  For each
  output slab $l$ and Cartesian index $i$, this produces a
  scalar.  The total number of scalars per
  $(\tau_m,\vq_n)$-point is $R^2 \times C \times N_\Sigma$,
  where $N_\Sigma$ is the number of output slab pairs within
  the self-energy bandwidth.
\end{enumerate}

An identical two-stage procedure produces $B^{j}_{l',rs}$
from $g_2(\tau_m;\vq_n)$, $\hat{h}_r(\vq_n)$, and
$\hat{f}'_s(\vq_n)$. No communication is required in this phase, each rank
operates only on its local $(\tau,\vq)$-points.

\paragraph{Phase~3: $\vq$-transpose
  (all-to-all within $\vq$-communicators).}
%
The FFT convolution in \cref{eq:sigma_sep} requires, for each
$(r,s,\tau_m)$, the values $A_{rs}(\tau_m;\vq_n)$ and
$B_{rs}(\tau_m;\vq_n)$ at \emph{all} $N_q$ wavevectors.
After Phase~2, these are distributed across $P_q$ ranks (each
holding $N_q^{\mathrm{loc}}$ wavevectors).  A
$\vq$-transpose---an all-to-all exchange within each
$\vq$-communicator---redistributes the data so that each rank
holds all $N_q$ wavevectors for a subset of
$(\tau,r,s,i,l)$-combinations.

The Communication volume per rank is
$N_\tau^{\mathrm{loc}} \times R^2 \times C \times N_\Sigma
  \times N_q^{\mathrm{loc}}$ complex numbers for $A$, and the
same for $B$.  Because of the compression in Phase~2, this is
a factor $M^2 N_B / (R^2 C N_\Sigma)$ smaller than the
hypothetical cost of transposing the full Green's function
matrices.

Without the separable approximation, $\vq$-distribution of the
$\vq$-space formulation \cref{eq:bubble_q} would require an
all-gather of $g(\tau;\vq')$ for all $\vq'$ at each $(\tau,\vq)$
working point, with communication volume
$N_\tau \times N_q \times M^2 N_B$ per rank---effectively
negating the memory saving from $\vq$-distribution.

\paragraph{Phase~4: $\vq$-FFT convolution
  (local after transpose).}
%
After the $\vq$-transpose, each rank holds $A^i_{l,rs}$ and
$B^j_{l',rs}$ at all $N_q$ wavevectors for its assigned subset
of $(r,s,\tau,i,l)$ and $(r,s,\tau,j,l')$ indices.  For each
such index combination:

\begin{enumerate}
\item Apply the inverse 2D FFT $\vq\to\vR$ to obtain
  $a^i_{l,rs}(\tau_m;\vR)$ and $b^j_{l',rs}(\tau_m;\vR)$ at all
  $N_q$ transverse lattice points $\vR$.
\item For each $\vR$: compute the outer product
  $c^{ij}_{ll',rs}(\tau_m;\vR) =
    a^i_{l,rs}(\tau_m;\vR) \cdot b^j_{l',rs}(\tau_m;\vR)$.
\item Accumulate over $(r,s)$ into the bubble:
  $\mathcal{B}^{ij}_{ll'}(\tau_m;\vR)
    \mathrel{+}= c^{ij}_{ll',rs}(\tau_m;\vR)$.
\item After all $(r,s)$ pairs: apply the forward 2D FFT
  $\vR\to\vq$ to recover
  $\mathcal{B}^{ij}_{ll'}(\tau_m;\vq)$.
\end{enumerate}

Alternatively, if one processes all $(r,s)$ pairs before the
final FFT, the accumulation is done in $\vR$-space (step~3)
and only one forward FFT per $(\tau,i,l,j,l')$ is needed
(step~4).  This reduces the number of forward FFTs from
$R^2$ to $1$ per index combination.

Cost per $\tau$-point is
$2 R^2 C N_\Sigma$ inverse FFTs (for $A$ and $B$) plus
$C^2 N_\Sigma^2$ forward FFTs (for $\mathcal{B}$), each of
length $N_q$.  With $N_q = 400$, $R = 10$, $C = 3$,
$N_\Sigma = 5$: approximately $3{,}000$ IFFTs and $225$ FFTs
per $\tau$-point.

\paragraph{Phase~5: reverse transposes and $\tau\to\omega$
  transform.}
%
The bubble $\mathcal{B}^{ij}_{ll'}(\tau_m;\vq_n)$ must be
transformed back to the $(\omega,\vq)$-distributed format for
injection into the next RGF solve.  This requires two steps:

\begin{enumerate}
\item \emph{Reverse $\vq$-transpose}: redistribute
  $\mathcal{B}(\tau;\vq)$ so that each rank holds all $\tau$
  at its local $\vq$-points, using an all-to-all within
  $\vq$-communicators.  The communicated data is the full
  self-energy matrix (not the compressed $A_{rs}$ scalars),
  so the volume is
  $N_\tau^{\mathrm{loc}} \times N_q^{\mathrm{loc}} \times
    C^2 N_\Sigma^2$.
\item \emph{$\tau\to\omega$ FFT}: at each $\vq_n$, apply the
  forward FFT $\tau\to\omega$ to obtain
  $\Sig^{ij}_{ll'}(\omega_k;\vq_n)$.
\item \emph{Reverse $\omega$-transpose}: redistribute from the
  NNZ layout (all $\omega$ at local entries) back to the
  Stacked layout (local $\omega$ at all entries) for the RGF
  solver.
\end{enumerate}

% -------------------------------------------------------
\subsubsection{Comparison to the non-separable parallelisation}
\label{sssub:compare_dist}
% -------------------------------------------------------

In the non-separable $\vq$-space formulation, the inner
$\vq'$-loop at each $(\tau,\vq)$-point accesses $G(\tau;\vq')$
for all $N_q$ values of $\vq'$.  If $G$ is distributed over
$P_q$ ranks, this requires either:
\begin{itemize}
\item an all-gather of $G(\tau;\cdot)$ at every $\tau$-step,
  with volume $N_q \times M^2 N_B$ per step, or
\item replication of $G(\tau;\vq)$ across all $\vq$-ranks,
  eliminating the memory benefit of $\vq$-distribution.
\end{itemize}
Neither option is viable for large $N_q$ and $M$.

The Guo et al.\ approach~\cite{guo2020} avoids this issue by
assigning one $\vq$-point per rank (full parallelisation of
$\vq$) and exchanging Green's function data between all pairs
of ranks at each SCBA iteration.  This works for small systems
($N_{BS} = 24$ in their silicon thin film, $N_q = 36$) but does
not scale to the nanoribbon regime ($N_{BS} = 3{,}408$,
$N_q = 400$) because the all-to-all volume grows as
$N_q^2 \times N_{BS}^2$.

The separable formulation resolves this by compressing
$M^2 N_B$ matrix elements to $R^2 C N_\Sigma$ scalars
\emph{before} the $\vq$-transpose.  The all-to-all volume in
Phase~3 is independent of $M$ and $N_B$, scaling only as
$R^2 C N_\Sigma \times N_\tau \times N_q$---a reduction by
the compression factor $M^2 N_B / (R^2 C N_\Sigma)$.

% -------------------------------------------------------
\subsubsection{Memory layout per rank}
\label{sssub:memory}
% -------------------------------------------------------

The peak memory per rank occurs during Phase~2 (bilinear
contraction), when the rank simultaneously holds:
\begin{center}
\begin{tabular}{lll}
  \toprule
  Array & Description & Size (complex numbers) \\
  \midrule
  $g(\tau;\vq)$
    & Green's function at current $\tau$,
    & $N_q^{\mathrm{loc}} \times N_{\mathrm{nnz}}^{\mathrm{loc}}$ \\
    & local $\vq$-points & \\[4pt]
  $\hat{f}_r(\vq), \hat{h}_r(\vq)$
    & Rank factors (replicated)
    & $2R \times N_q^{\mathrm{loc}} \times M\hat{N}_c$ \\[4pt]
  $\vv_s(\tau;\vq)$
    & Intermediate vectors
    & $R \times N_q^{\mathrm{loc}} \times M\hat{N}_c$ \\[4pt]
  $A_{rs}(\tau;\vq)$, $B_{rs}(\tau;\vq)$
    & Compressed scalars
    & $2R^2 \times N_q^{\mathrm{loc}} \times C N_\Sigma$ \\[4pt]
  $\mathcal{B}(\tau;\vq)$
    & Accumulated bubble (all $\tau$)
    & $N_\tau \times N_q^{\mathrm{loc}} \times C^2 N_\Sigma^2$ \\
  \bottomrule
\end{tabular}
\end{center}
The bubble $\mathcal{B}$ must persist across all $\tau$-steps
(to perform the final $\tau\to\omega$ FFT), while $g$,
$\vv_s$, $A_{rs}$, and $B_{rs}$ are overwritten at each
$\tau$-step.  In the separable formulation, $g$ needs to be
held only for the current $\tau$-step, not for all $\tau$
simultaneously---in contrast to the non-separable $\vq$-space
method, which requires random access to $g(\tau;\vq')$ and
$g(\tau;\vq-\vq')$ across the full $\tau$-range.

% -------------------------------------------------------
\subsubsection{Communication summary}
\label{sssub:comm_summary}
% -------------------------------------------------------

\Cref{tab:comm} summarises the four all-to-all exchanges per
SCBA iteration.  The $\omega$-transposes (Phases~1 and~5c) are
identical to those in the existing QuaTrEx electron--electron
self-energy computation.  The $\vq$-transposes (Phases~3
and~5a) are new and operate on the compressed $A_{rs}$,
$B_{rs}$ scalars or the self-energy matrix $\mathcal{B}$
respectively.

\begin{center}
\begin{tabular}{llll}
  \toprule
  Phase & Communicator & Direction & Volume per rank \\
  \midrule
  1 ($\omega\to\tau$)
    & $\omega$-comm
    & Stacked$\to$NNZ
    & $N_\omega^{\mathrm{loc}} N_q^{\mathrm{loc}}
      N_{\mathrm{nnz}}^{\mathrm{loc}}$ \\[3pt]
  3 ($\vq$-transpose)
    & $\vq$-comm
    & local-$\vq\to$all-$\vq$
    & $N_\tau N_q^{\mathrm{loc}} R^2 C N_\Sigma$ \\[3pt]
  5a (reverse $\vq$)
    & $\vq$-comm
    & all-$\vq\to$local-$\vq$
    & $N_\tau N_q^{\mathrm{loc}} C^2 N_\Sigma^2$ \\[3pt]
  5c ($\tau\to\omega$)
    & $\omega$-comm
    & NNZ$\to$Stacked
    & $N_\omega^{\mathrm{loc}} N_q^{\mathrm{loc}}
      N_{\mathrm{nnz}}^{\mathrm{loc}}$ \\
  \bottomrule
\end{tabular}
\label{tab:comm}
\end{center}

The $\vq$-transpose in Phase~3 is the communication
bottleneck specific to the phonon--phonon self-energy.  Its
volume scales as $R^2 C N_\Sigma$ per $(\tau,\vq)$-point,
which is independent of the cross-section size $M$.  For
$R = 10$, $C = 3$, $N_\Sigma = 5$: $1{,}500$ complex numbers
per $(\tau,\vq)$-point, compared to
$M^2 N_B \approx 2.2\times 10^6$ for the full Green's
function---a reduction by three orders of magnitude.

% -------------------------------------------------------
\subsubsection{Computational cost per rank per SCBA iteration}
\label{sssub:cost_per_rank}
% -------------------------------------------------------

The dominant cost is Phase~2 (bilinear contraction).  Per
rank, the number of $(\tau,\vq)$-points processed is
$N_\tau \times N_q^{\mathrm{loc}}$.  At each point:

\begin{itemize}
\item Step~1 (matrix--vector products): $R$ products, each
  involving $\hat{N}_c$ block-rows of the block-tridiagonal $g$
  at cost $\cO(\hat{N}_c N_{BS}^2)$.
  Total: $R\,\hat{N}_c\,N_{BS}^2$ per $(\tau,\vq)$.
\item Step~2 (inner products): $R^2$ inner products of length
  $M\hat{N}_c$, for each of $C \times N_\Sigma$ output pairs.
  Total: $R^2\,M\hat{N}_c\,C\,N_\Sigma$ per $(\tau,\vq)$.
\end{itemize}

For $P$ and $Q$ together, the per-rank cost is
\begin{equation}
  \cW_{\mathrm{rank}} =
  N_\tau \frac{N_q}{P_q}
  \bigl[
    2R\,\hat{N}_c\,N_{BS}^2
  + 2R^2\,M\hat{N}_c\,C\,N_\Sigma
  \bigr]
  / P_S.
  \label{eq:cost_per_rank}
\end{equation}
The first term dominates for the nanoribbon parameters of
\cref{tab:nanoribbon}: with $N_{BS} = 3{,}408$,
$\hat{N}_c = 4$, $R = 10$, $N_\tau = 500$,
$N_q^{\mathrm{loc}} = N_q/P_q$, each rank performs
$\cO(10^{14}/P_q P_S)$ floating-point operations per SCBA
iteration.


% ============================================================
\subsection{No Tranverse Periodicity}
\label{sub:nq1}
% ============================================================

The analysis above assumed a nonzero transverse wavevector
grid ($N_q > 1$), which arises when the device cross-section
has periodic boundary conditions in one or both transverse
directions.  For narrow nanoribbons, FinFETs, or gate-all-around
nanowire transistors whose cross-section is fully confined
(no transverse periodicity), $N_q = 1$ and the self-energy
\cref{eq:sigma_phph} reduces to the purely 1D form of
Guo et~al.~\cite[Eq. 8]{guo_quantum_2020}:
\begin{equation}
  \Sig^{\gtrless,ij}_{ll'}(\omega)
  = \frac{i\hbar}{2}
    \sum_{\substack{l_{1..4}\\j_{1..4}}}
    \int\!\frac{d\omega'}{2\pi}\;
    \Phi^{ij_1j_2}_{ll_1l_2}\;
    \Phi^{j_3j_4 j}_{l'l_3l_4}\;
    G^{\gtrless,j_1j_4}_{l_1 l_4}(\omega')\;
    G^{\gtrless,j_2j_3}_{l_2 l_3}(\omega{-}\omega').
  \label{eq:sigma_1d}
\end{equation}
The total cost per SCBA iteration is then
\begin{align}
  \cW_q\big|_{N_q=1}
  &= N_\tau\;\cC_\mathrm{ic}
  = N_\tau\,M^3\,\hat{N}_c^2\,N_B,
  \tag{non-separable, $N_q{=}1$}
  \label{eq:cost_nq1_nonsep}
  \\[4pt]
  \cW_S\big|_{N_q=1}
  &= N_\tau\bigl[
    R\,M^2\hat{N}_c N_B + R^2 M\hat{N}_c
  \bigr].
  \tag{separable, $N_q{=}1$}
  \label{eq:cost_nq1_sep}
\end{align}
The separable formulation still provides a speedup of
\begin{equation}
  \frac{\cW_q}{\cW_S}\bigg|_{N_q=1}
  \approx
  \frac{M\,\hat{N}_c}{R},
  \label{eq:speedup_nq1}
\end{equation}
from the bilinear-form reduction ($M^3 \to M^2$) and shared-$G$
reuse (factor~$R$).

The distributed algorithm simplifies accordingly: only
the $\omega$-distribution and spatial decomposition from
QuaTrEx are needed.  Phase~1 (the existing
Stacked$\,\to\,$NNZ transpose for $\omega\to\tau$) and the
bilinear contraction (Phase~2 at a single $\vq$-point) proceed
as before.  Phases~3--4 (the $\vq$-transpose and FFT
convolution) are absent.  Phase~5 reduces to the
$\tau\to\omega$ FFT and the reverse $\omega$-transpose.

